%% SIAM Journal on Scientific Computing -- prepared with siamart220329.cls
%% Requires: \documentclass[review,hidelinks]{siamart220329}
%% Using article as stand-in until siamart220329.cls is downloaded
\documentclass[12pt,a4paper]{article}


\usepackage{graphicx}
\usepackage{amsmath,amsfonts,amssymb}
\usepackage{booktabs}
\usepackage{array}
\usepackage{multirow}
\usepackage{xcolor}
\usepackage{algorithm}
\usepackage{algpseudocode}
\usepackage{url}
\usepackage{float}
\usepackage[caption=false]{subfig}
\usepackage{microtype}
\usepackage{enumitem}
\usepackage{boldline}
\usepackage{afterpage}
\usepackage{amsthm}
\usepackage[margin=1in]{geometry}
\usepackage{setspace}
\doublespacing

% SIAM theorem environments (use amsthm defs already present)


% Theorem environments for SIAM-style formal claims
\newtheorem{proposition}{Proposition}
\newtheorem{theorem}{Theorem}
\newtheorem{remark}{Remark}

\graphicspath{{plots/}}

% ─────────────────────────────────────────────────────────────────
%  TITLE
% ─────────────────────────────────────────────────────────────────
\title{A Comparative Study of Hamiltonian Hyperparameter Dynamics,
       Hybrid Optimisation, and a Unified Framework for
       Self-Tuning Hyperparameter Optimization}

\author{Anonymous Author\thanks{Submitted to SIAM Journal on Scientific Computing (SISC).}}

% ─────────────────────────────────────────────────────────────────
\begin{document}
\maketitle
% SIAM running head: \headers{Hamiltonian Hyperparameter Dynamics}{Author}

% ══════════════════════════════════════════════════════════════════
%  PAGE 1 — ABSTRACT + NOMENCLATURE
% ══════════════════════════════════════════════════════════════════

\begin{abstract}
Hyperparameter optimisation (HPO) is a fundamental bottleneck in
modern machine learning systems, often dictating both convergence
behaviour and final model performance.  Classical approaches such as
grid search and random search suffer from exponential scaling, while
Bayesian optimisation improves sample efficiency but remains
computationally expensive and inherently decoupled from training
dynamics.  In parallel, advances in optimisation algorithms—ranging
from adaptive first-order methods such as Adam to quasi-Newton methods
such as L-BFGS—have improved convergence rates but do not directly
address hyperparameter evolution during training.
In this work, we propose a unified perspective that models optimisation
as a continuous dynamical system governed by Hamiltonian mechanics.  We
present three complementary frameworks: \textbf{(A) Hamiltonian
Hyperparameter Dynamics (HHD)}, which formulates joint optimisation of
weights and hyperparameters via Hamiltonian mechanics and symplectic
integration; \textbf{(B) a Hybrid Adam $+$ L-BFGS $+$ Bayesian
Optimisation (ABBO)} framework representing a strong, practically
competitive baseline; and \textbf{(C) a novel Unified HHD-ABBO
framework} that integrates stochastic gradient descent, symplectic
sampling, and curvature-based refinement within a single, coherent
training curriculum applied at every epoch.
We evaluate these approaches across four benchmark categories: a
controlled physics benchmark (harmonic oscillator energy landscape),
a real-world CNN classification task (MNIST), and 11 standardised
tabular benchmarks from HPOBench, HPOLib, and NAS-Bench-201.  Results
demonstrate that the Unified framework (Method~C) achieves the
strongest reconstruction accuracy on the physics benchmark
($\mathrm{MSE} = 0.0033$, $R^2 = 0.9999$), outperforming ABBO by
$30\times$ in error magnitude while retaining symplectic conservation,
smooth hyperparameter trajectories, and second-order curvature
exploitation.  On standardised HPO benchmarks, Method~C attains the
best average rank ($2.18$) across all 11 datasets, outperforming
Optuna TPE ($2.64$), Random Search ($3.55$), and both individual
frameworks.  The Unified method achieves $97.70\,\%$ CNN validation
accuracy in $3.3\times$ less wall time than the BO-based baseline.
This work establishes a principled bridge between physics-based
optimisation and modern deep learning, enabling structured,
interpretable, and efficient learning systems.
\end{abstract}

\begin{keywords}
Hamiltonian mechanics, hyperparameter optimisation, symplectic integration,
Bayesian optimisation, meta-learning, Hamiltonian Monte Carlo, neural networks
\end{keywords}

\begin{MSCcodes}
68T05, 65P10, 49M37, 62C10
\end{MSCcodes}

% ─────────────────────────────────────────────────────────────────
\section{Nomenclature}

\begin{table}[tb]
\centering
\caption{Complete Symbol Reference}
\label{tab:nomenclature}
\renewcommand{\arraystretch}{1.35}
\begin{tabular}{>{\bfseries}p{1.45cm} p{6.0cm}}
\toprule
Symbol & Description \\
\midrule
$q$              & Generalised position (phase-space coordinate) \\
$p$              & Generalised momentum \\
$\theta$         & Network weight vector $\in\mathbb{R}^d$ \\
$\lambda$        & Hyperparameter vector $\in\mathbb{R}^k$ \\
$p_\theta$       & Conjugate momentum for network weights \\
$p_\lambda$      & Conjugate momentum for hyperparameters \\
$m_\theta$       & Effective mass (inertia) of weight momenta \\
$m_\lambda$      & Effective mass (inertia) of HP momenta \\
$H$              & Total Hamiltonian $= T + V$ \\
$T$              & Kinetic energy $= \sum p^2 / (2m)$ \\
$V = \mathcal{L}$& Potential energy $=$ training loss \\
$\mathcal{L}$    & Mean-Squared Error (MSE) loss function \\
$\varepsilon$    & Leapfrog step size \\
$L$              & Number of leapfrog sub-steps per HMC proposal \\
$\eta,\alpha$    & Learning rates (Adam / L-BFGS line search) \\
$\beta_1,\beta_2$& Adam first- and second-moment decay rates \\
$m_t, v_t$       & Adam running first and second moment estimates \\
$\hat{m}_t,\hat{v}_t$ & Bias-corrected Adam moment estimates \\
$H_t^{-1}$      & L-BFGS implicit inverse Hessian approximation \\
$k(\cdot,\cdot)$ & Gaussian Process kernel function (RBF) \\
$\mu(\lambda)$   & GP posterior predictive mean at $\lambda$ \\
$\sigma(\lambda)$& GP posterior predictive standard deviation \\
$\mathrm{EI}(\lambda)$ & Expected Improvement acquisition function \\
$\Phi,\phi$      & Standard normal CDF and PDF \\
$f^*$            & Current best observed validation loss \\
$\xi$            & Exploration–exploitation trade-off parameter \\
$N$              & Number of training samples \\
$d$              & Dimensionality of weight space \\
$k$              & Dimensionality of hyperparameter space \\
$\sigma_f^2$     & GP signal variance \\
$\ell$           & GP length-scale parameter \\
$\sigma_n^2$     & GP observation noise variance \\
MAE              & Mean Absolute Error on the energy landscape \\
RMSE             & Root Mean Squared Error on the energy landscape \\
$R^2$            & Coefficient of determination \\
$\Delta H$       & Hamiltonian error of a leapfrog proposal \\
$P$              & Plateau-detection patience window (epochs) \\
$\tau$           & Plateau-detection relative-loss tolerance \\
\bottomrule
\end{tabular}
\end{table}

\newpage

% ══════════════════════════════════════════════════════════════════
%  PAGE 2 — INTRODUCTION (dedicated page)
% ══════════════════════════════════════════════════════════════════

\section{Introduction}

\subsection{The Genesis and Evolution of Hyperparameter Optimisation}

Hyperparameter Optimisation (HPO) finds its origins in the early days
of statistical modelling, where parameters such as regularisation
coefficients were tuned via manual experimentation or exhaustive Grid
Search~\cite{lecun2015,raiaan2024}.  The formalisation of HPO as a
distinct field gained decisive momentum with the landmark study of
Bergstra and Bengio (2012)~\cite{bergstra2012}, who demonstrated that
Random Search was significantly more efficient than Grid Search by
exploiting the empirically observed \emph{low effective dimensionality}
of most hyperparameter landscapes.  Their analysis showed that, for a
fixed evaluation budget, unstructured random sampling tends to cover
more of the truly influential dimensions than an exhaustive grid, which
wastes evaluations on directions that do not appreciably affect
generalisation.

As deep architectures grew in complexity, the black-box nature of
neural network training necessitated strategies that were both
sample-efficient and capable of incorporating prior knowledge from
previous evaluations.  This requirement led to the widespread adoption
of Sequential Model-Based Optimisation (SMBO) and, in particular,
Bayesian Optimisation (BO)~\cite{snoek2012}.  By fitting a Gaussian
Process (GP) surrogate model over the objective function and choosing
new candidate hyperparameter configurations through an acquisition
function—most commonly Expected Improvement (EI)~\cite{mockus1975}—BO
can direct the search towards promising regions of the hyperparameter
space while balancing exploration and exploitation.  Despite its
sample efficiency, BO remains computationally expensive in the
outer-loop sense: each trial requires a complete re-initialisation and
re-training of the underlying model, and the GP itself scales
cubically with the number of observations.

\subsection{The Hamiltonian Tradition: From Mechanics to Machine Learning}

The dynamical framework employed in this paper traces its lineage to
Sir William Rowan Hamilton's 1833 reformulation of Newtonian mechanics.
Rather than describing motion in terms of forces and accelerations,
Hamiltonian mechanics shifts focus to \emph{energy conservation} and
\emph{canonical coordinates} $(q, p)$, where $q$ denotes generalised
positions and $p$ their conjugate momenta.  This framework proved
foundational for celestial mechanics and, later, for quantum mechanics.

In the late twentieth century, Hamiltonian dynamics crossed the
boundary from pure physics into computational statistics.  Duane
\emph{et al.}~\cite{duane1987} introduced Hybrid Monte Carlo (HMC)—now
universally called Hamiltonian Monte Carlo—exploiting the symplectic
structure of the leapfrog integrator to allow a Markov Chain to make
large, informed proposals across a high-dimensional probability
distribution without the diffusive random-walk behaviour that afflicts
Metropolis-Hastings algorithms.  Neal's comprehensive
treatment~\cite{neal2011} subsequently established HMC as a reference
method for Bayesian posterior sampling, while Betancourt's conceptual
introduction~\cite{betancourt2017} articulated why symplecticity is
precisely the geometrical property that gives HMC its efficiency: the
leapfrog integrator conserves a \emph{shadow Hamiltonian} close to the
true one, bounding numerical energy error globally and ensuring high
Metropolis acceptance rates.

\subsection{Convergent Motivation: Treating Hyperparameters as Particles}

Despite their individual successes, Bayesian HPO and Hamiltonian
sampling have historically remained decoupled components of the machine
learning workflow.  Bayesian HPO treats the model as a static black
box, and the surrogate GP is updated only between complete training
runs.  No information about the trajectory of weights during training
is incorporated into the hyperparameter update.  Similarly, standard
HMC is used for posterior inference over weights or model parameters,
not for the simultaneous, continuous evolution of hyperparameters.

Our central motivation arises from a unifying question: \emph{Can we
treat hyperparameters as particles subject to the same dynamical laws as
network weights?}  If the training loss defines a potential energy
surface $V = \mathcal{L}(\theta, \lambda)$ over the joint space of
weights and hyperparameters, then introducing conjugate momenta
$(p_\theta, p_\lambda)$ allows both to evolve together under Hamilton's
equations.  This perspective eliminates the wasteful ``outer trial
loop'' of BO, moves hyperparameter tuning inside the training loop, and
naturally inherits the symplectic conservation properties of HMC.

Crucially, physics-inspired neural network learning~\cite{greydanus2019,zhang2025}
has demonstrated that networks trained on physical systems benefit
enormously from optimisers that respect the geometric structure of the
underlying dynamics.  The harmonic oscillator testbed chosen in this
work amplifies these benefits: because the ground-truth Hamiltonian is
known, all metrics—MAE, RMSE, $R^2$—can be evaluated exactly, making it
an ideal benchmark for comparing the fidelity of competing frameworks.

\subsection{Contributions}

This paper makes the following contributions:

\begin{enumerate}[leftmargin=*, label=\arabic*.]
\item We formulate \textbf{Hamiltonian Hyperparameter Dynamics (HHD)}
      (Method A), a principled framework in which network weights
      \emph{and} hyperparameters are jointly evolved via symplectic
      leapfrog integration over an augmented Hamiltonian.

\item We construct a \textbf{Hybrid ABBO baseline} (Method B) combining
      Adam, L-BFGS, and Bayesian Optimisation as the strongest
      practically competitive reference.

\item We introduce the \textbf{Unified HHD-ABBO framework}
      (Method C), which integrates all three optimisers into a
      single three-phase curriculum applied within every training
      epoch, inheriting symplectic conservation from A and
      second-order curvature from B while eliminating the expensive
      outer BO loop.

\item We provide a rigorous \textbf{multi-metric evaluation} on the
      harmonic oscillator energy landscape (MAE, RMSE, $R^2$,
      Hamiltonian conservation, smooth HP trajectories, wall time),
      on \textbf{CNN classification on MNIST}, and across
      \textbf{11 standardised tabular benchmarks} from HPOBench,
      HPOLib, and NAS-Bench-201~\cite{eggensperger2021,dong2020nasbench201}.
\end{enumerate}

\subsection{Research Questions}

\begin{enumerate}[leftmargin=*, label=\textbf{RQ\arabic*.}]
\item Can network weights and hyperparameters be jointly optimised via
      Hamiltonian dynamics, and does this preserve physical properties
      such as energy conservation?
\item How does the pure HHD approach compare quantitatively and
      qualitatively to the practically competitive ABBO baseline?
\item Does the Unified method, by combining all three optimisers in a
      structured curriculum, outperform both individual frameworks?
\item What unique theoretical or practical properties does each
      framework provide, and in which deployment scenarios should
      each be preferred?
\end{enumerate}



% ══════════════════════════════════════════════════════════════════
%  SECTION 2 — RELATED WORK
% ══════════════════════════════════════════════════════════════════

\section{Related Work}
\label{sec:related}

\subsection{Gradient-Based Hyperparameter Optimisation}

Bilevel optimisation methods treat HPO as an outer minimisation over
validation loss subject to an inner minimisation of training loss.  The
Lorraine et al.\ approach~\cite{lorraine2020} approximates the
hypergradient through the implicit function theorem, avoiding
unrolled differentiation at the cost of a linear solve per step.
Franceschi et al.~\cite{franceschi2018} study the forward-mode
bilevel variant and show that a $T$-step unrolled gradient is a
biased estimate of the true hypergradient.  MAML~\cite{finn2017maml}
adapts model-agnostic meta-learning to tune initialisation, which
can be interpreted as a bilevel problem over the task distribution.
DARTS~\cite{liu2019darts} relaxes the discrete architecture search
problem via a continuous relaxation, enabling gradient-based
optimisation of architecture hyperparameters jointly with weights.
By contrast, our Method A and Method C avoid bilevel structure
entirely by treating hyperparameters as particles in a joint
Hamiltonian system, eliminating both the inner loop and the
hypergradient computation.

\subsection{Physics-Informed Optimisers and Symplectic Methods}

Hamiltonian Neural Networks (HNNs)~\cite{greydanus2019} learn
conservation laws directly from data by parameterising the Hamiltonian
and ensuring that predictions respect Hamilton's equations.  Lagrangian
Neural Networks~\cite{cranmer2020} take an analogous approach via the
Euler--Lagrange equations.  On the optimisation side, Franca et
al.~\cite{franca2020} establish connections between symplectic
integrators and stochastic gradient descent, showing that variants of
Adam can be derived from a discrete Hamiltonian.  Zhang et
al.~\cite{zhang2025} apply Hamiltonian dynamics to physics-informed
nets for energy landscape learning—the same testbed we use in this
work.  Our contribution extends this line of research by using
Hamiltonian dynamics not to learn a physical system's equations, but
to govern the optimisation process itself.

\subsection{Multi-Fidelity Hyperparameter Optimisation}

Hyperband~\cite{li2018hyperband} treats early-stopping as a pure
exploration problem and applies successive halving to a random
portfolio of configurations, achieving near-optimal sample complexity
without a surrogate model.  BOHB~\cite{falkner2018bohb} combines
Hyperband's successive-halving scheduling with SMAC's KDE-based
Bayesian optimisation~\cite{bergstra2011tpe} to acquire configurations
that are both sample-efficient and budget-aware.  ASHA~\cite{li2020asha}
extends Hyperband to the asynchronous parallel setting.  While these
methods are strong baselines on tabular benchmarks—and BOHB and SMAC3
are included in our Phase 3 comparison—they operate in a conventional
outer-loop paradigm and do not preserve any physical structure of the
loss landscape.

\subsection{Neural Architecture Search}

NAS-Bench-201~\cite{dong2020nasbench201} provides a tabular benchmark
of architecture performance that enables reproducible comparison of
search strategies without training overhead.  Regularised
evolution~\cite{real2019regularized} applies population-based
evolutionary search to architecture spaces.  Our work does not focus
on NAS, but we include NAS-Bench-201 in our benchmark suite to
demonstrate that Method C's continuous HP relaxation generalises
beyond regression and tabular classification tasks.

\section{Background and Theoretical Foundations}

\subsection{Adam Optimiser}

Adam~\cite{kingma2015adam} maintains per-parameter adaptive learning
rates through exponential moving averages of the gradient $g_t$ and
its element-wise square.  Given gradient $g_t$ at step $t$, the
update rules are:
\begin{align}
m_t &= \beta_1 m_{t-1} + (1-\beta_1)g_t \label{eq:adam1}\\
v_t &= \beta_2 v_{t-1} + (1-\beta_2)g_t^2 \label{eq:adam2}\\
\theta_{t+1} &= \theta_t - \eta\,\frac{\hat{m}_t}{\sqrt{\hat{v}_t}+\epsilon}
\label{eq:adam3}
\end{align}
where the bias-corrected estimates are $\hat{m}_t = m_t/(1-\beta_1^t)$
and $\hat{v}_t = v_t/(1-\beta_2^t)$.  Default hyperparameters are
$\beta_1 = 0.9$, $\beta_2 = 0.999$, and $\epsilon = 10^{-8}$.

Adam excels in the early training regime where the loss landscape is
far from any stationary point: its adaptive step sizes accommodate
widely varying gradient magnitudes across parameters, and the moment
estimates provide implicit momentum that accelerates traversal of
shallow ravines.  Its per-step computational cost is $\mathcal{O}(d)$,
making it well suited for mini-batch stochastic settings.  However,
Adam does not incorporate second-order curvature information, which
limits its convergence rate near optima.

\subsection{L-BFGS (Limited-Memory Broyden--Fletcher--Goldfarb--Shanno)}

L-BFGS~\cite{liu1989lbfgs,byrd1995} approximates Newton's method for
large-scale problems by implicitly representing the inverse Hessian
$H_t^{-1}$ through a rolling window of the $m$ most recent
gradient-difference vector pairs $\{(s_i, y_i)\}_{i=t-m}^{t-1}$,
where $s_i = \theta_i - \theta_{i-1}$ and
$y_i = \nabla\mathcal{L}(\theta_i) - \nabla\mathcal{L}(\theta_{i-1})$.
The parameter update is:
\begin{equation}
\theta_{t+1} = \theta_t - \alpha_t\,H_t^{-1}\,g_t
\label{eq:lbfgs}
\end{equation}
where the step size $\alpha_t$ is determined by a backtracking line
search satisfying the strong Wolfe conditions~\cite{wolfe1969}.  The
two-loop recursion for the L-BFGS direction requires only
$\mathcal{O}(md)$ memory and $\mathcal{O}(md)$ computation per step,
avoiding the $\mathcal{O}(d^2)$ storage of the full Hessian.  History
size $m = 10$ is used throughout this work.

Near stationary points, L-BFGS exhibits super-linear convergence,
making it the method of choice for full-batch fine-tuning once the
Adam warm-up has located a promising basin.  Its reliance on the full
gradient (or a large-batch gradient) makes it unsuitable for the
early, stochastic phase of training where gradient estimates are noisy.

\subsection{Bayesian Optimisation with Gaussian Process Surrogates}

Bayesian Optimisation~\cite{snoek2012} models the black-box objective
$f: \lambda \mapsto \mathcal{L}_{val}$ with a Gaussian Process
prior~\cite{rasmussen2006}.  The RBF (squared exponential) kernel
encodes the prior belief that nearby hyperparameter configurations
produce similar validation losses:
\begin{equation}
k(\lambda,\lambda') = \sigma_f^2\exp\!\left(-\frac{\|\lambda-\lambda'\|^2}{2\ell^2}\right)
+\sigma_n^2\delta_{\lambda\lambda'}
\label{eq:rbf}
\end{equation}
After $t$ observations $\mathcal{D}_t = \{(\lambda_i, f_i)\}_{i=1}^t$,
the GP provides a closed-form posterior predictive distribution
$\mathcal{N}(\mu_t(\lambda), \sigma_t^2(\lambda))$ over any
unobserved $\lambda$.  New candidates are selected by maximising the
Expected Improvement acquisition function~\cite{mockus1975}:
\begin{align}
\mathrm{EI}(\lambda) &=
  \bigl(f^*-\mu(\lambda)-\xi\bigr)\Phi(Z)
  +\sigma(\lambda)\phi(Z)
\label{eq:ei}\\
Z &= \frac{f^*-\mu(\lambda)-\xi}{\sigma(\lambda)} \notag
\end{align}
where $f^*$ is the current best observed validation loss,
$\xi \ge 0$ is an exploration bonus, and $\Phi, \phi$ are the
standard normal CDF and PDF respectively.  The acquisition surface is
maximised via multi-start L-BFGS-B, with random restarts to mitigate
local optima.

The GP accumulates knowledge across all trials through its posterior,
and each successive query is directed by the full history of
observations.  Despite this sample efficiency, BO incurs a fixed cost
of one complete model training run per trial and GP inference that
scales cubically with the number of completed trials.

\subsection{Hamiltonian Monte Carlo}

HMC~\cite{duane1987,neal2011} augments the target parameter space with
auxiliary momentum variables $p \sim \mathcal{N}(0, M)$ and defines
an extended energy:
\begin{equation}
H(\theta,p) = \mathcal{L}(\theta) + \tfrac{1}{2}p^TM^{-1}p
\label{eq:hmc}
\end{equation}
The system evolves under Hamilton's equations, numerically integrated
by the \emph{leapfrog} (St\"ormer--Verlet) integrator~\cite{leimkuhler2004}:
\begin{align}
p_{t+\varepsilon/2} &= p_t - \tfrac{\varepsilon}{2}\nabla_\theta\mathcal{L}(\theta_t)
\label{eq:lf1}\\
\theta_{t+\varepsilon} &= \theta_t + \varepsilon M^{-1}p_{t+\varepsilon/2}
\label{eq:lf2}\\
p_{t+\varepsilon}   &= p_{t+\varepsilon/2}
  - \tfrac{\varepsilon}{2}\nabla_\theta\mathcal{L}(\theta_{t+\varepsilon})
\label{eq:lf3}
\end{align}
After $L$ leapfrog steps, the proposed state $(\theta', p')$ is
accepted with probability $\min(1, e^{-\Delta H})$, where
$\Delta H = H(\theta', p') - H(\theta, p)$ is the Hamiltonian error
introduced by discretisation.  Because the leapfrog integrator is
time-reversible and volume-preserving (symplectic), it conserves a
\emph{shadow Hamiltonian} $\tilde{H} = H + \mathcal{O}(\varepsilon^2)$,
keeping $\Delta H$ bounded and ensuring high acceptance rates even for
large step sizes~\cite{hairer2006,betancourt2017}.

The mass matrix $M$ acts as a preconditioner: setting $M$ to the
diagonal of the Hessian at the current point produces an effective
rescaling of the parameter space that equalises curvature and improves
mixing.

\newpage

% ══════════════════════════════════════════════════════════════════
%  SECTION 4 — METHODOLOGY
% ══════════════════════════════════════════════════════════════════

\section{Methodology}

This section describes the mathematical formulations and design
decisions that differentiate the three frameworks.  Concrete
hyperparameter values and training settings are deferred to
Section~\ref{sec:setup}.

\subsection{Method A: Hamiltonian Hyperparameter Dynamics (HHD)}

\subsubsection{Augmented Phase Space}

The central insight of Method A is to treat both network weights
$\theta \in \mathbb{R}^d$ and hyperparameters
$\lambda \in \mathbb{R}^k$ as \emph{generalised positions} in a joint
phase space, equipping each with a conjugate momentum.  Following the
physics-informed learning framework of Greydanus
\emph{et al.}~\cite{greydanus2019} and Zhang \emph{et al.}~\cite{zhang2025},
we define the augmented Hamiltonian:
\begin{equation}
H(\theta,\lambda,p_\theta,p_\lambda)
  = \underbrace{\sum_i\frac{p_{\theta,i}^2}{2m_\theta}
    + \sum_j\frac{p_{\lambda,j}^2}{2m_\lambda}}_{T\;(\text{kinetic})}
  + \underbrace{\mathcal{L}(\theta,\lambda)}_{V\;(\text{potential})}
\label{eq:augH}
\end{equation}
with separate effective masses $m_\theta$ for weight momenta and
$m_\lambda$ for HP momenta.  The mass ratio $m_\lambda / m_\theta$
controls the relative \emph{timescale} of HP evolution: a smaller
$m_\lambda$ allows hyperparameters to move more rapidly relative to
weights, enabling exploration of the HP landscape without destabilising
weight convergence.

\subsubsection{Hamilton's Equations for the Joint System}

Differentiating Eq.~\eqref{eq:augH} with respect to the four
canonical variables yields four coupled ordinary differential
equations:
\begin{align}
\dot\theta      &= +\frac{\partial H}{\partial p_\theta}
                 = \frac{p_\theta}{m_\theta}
\label{eq:heq1}\\
\dot{p}_\theta  &= -\frac{\partial H}{\partial \theta}
                 = -\nabla_\theta\mathcal{L}(\theta,\lambda)
\label{eq:heq2}\\
\dot\lambda     &= +\frac{\partial H}{\partial p_\lambda}
                 = \frac{p_\lambda}{m_\lambda}
\label{eq:heq3}\\
\dot{p}_\lambda &= -\frac{\partial H}{\partial \lambda}
                 = -\nabla_\lambda\mathcal{L}(\theta,\lambda)
\label{eq:heq4}
\end{align}
Weight gradients $\nabla_\theta\mathcal{L}$ are computed via
automatic differentiation.  Because discrete hyperparameters (network
depth, width, batch size) are not amenable to autodiff, HP gradients
$\nabla_\lambda\mathcal{L}$ are approximated by central finite
differences~\cite{rathore2024}.  The joint leapfrog integrator applies
Eqs.~\eqref{eq:lf1}--\eqref{eq:lf3} simultaneously to all four
variables $(\theta, \lambda, p_\theta, p_\lambda)$.

\subsubsection{Continuous Relaxation of Discrete Hyperparameters}

To allow hyperparameters to evolve continuously within the Hamiltonian
framework, discrete variables are replaced by smooth surrogates as
shown in Table~\ref{tab:hpparam}.

\begin{table}[tb]
\centering
\caption{Continuous Parameterisations for Discrete Hyperparameters}
\label{tab:hpparam}
\renewcommand{\arraystretch}{1.2}
\begin{tabular}{llll}
\toprule
Hyperparameter & Variable & Parameterisation & Range \\
\midrule
Learning rate  & $\lambda_{lr}$ & $\log_{10}(\eta)$    & $(-4,-1)$ \\
Dropout rate   & $\lambda_{do}$ & $p_{\mathrm{drop}}$  & $(0,\,0.3)$ \\
Network depth  & $\lambda_{nl}$ & Continuous integer   & $(1,\,8)$ \\
Layer width    & $\lambda_{nn}$ & Continuous integer   & $(16,\,256)$ \\
Batch size     & $\lambda_{bs}$ & $\log_2(b)$          & $(4,\,6)$ \\
\bottomrule
\end{tabular}
\end{table}

The log-scale parameterisation of learning rate and batch size ensures
that the Hamiltonian dynamics operate on a space where unit displacements
correspond to multiplicative changes in the original variable—a more
natural geometry for these exponentially-varying quantities.  After
each HMC proposal, continuous values are snapped to the nearest valid
integer for architecture parameters before being used to instantiate
the network.

\subsubsection{Two-Phase Training Procedure}

Method A employs a two-phase training procedure.  In Phase 1, an Adam
warm-up trains the network weights for a fixed number of epochs without
any HP update.  This provides a sensible initialisation of the weight
landscape before committing to the Hamiltonian trajectory.  In Phase 2,
the full joint leapfrog-and-Metropolis--Hastings loop co-evolves
$(\theta, \lambda)$.  Adam micro-steps continue to be applied to the
weights after each accepted HMC proposal, providing the stochastic
gradient signal needed to navigate mini-batch noise while the
Hamiltonian dynamics explore the HP space.

\subsection{Method B: Hybrid Adam + L-BFGS + Bayesian Optimisation (ABBO)}

\subsubsection{Decoupled Outer--Inner Loop Structure}

The ABBO framework~\cite{snoek2012,liu1989lbfgs,kingma2015adam}
separates hyperparameter search (outer loop) from weight training
(inner loop).  The outer loop is a standard Bayesian Optimisation
procedure driven by a GP surrogate and the Expected Improvement
acquisition of Eq.~\eqref{eq:ei}.  For each hyperparameter
configuration $\lambda_t$ proposed by the acquisition function (or
sampled uniformly during the random initialisation phase), a fresh
model is constructed and trained to convergence in the inner loop.

\subsubsection{Inner Training Curriculum}

The inner training loop itself is a hybrid of two complementary
optimisers.  Adam is run first for $N_a$ epochs, using its adaptive
learning rates to make rapid progress from a random initialisation.
Once Adam has converged to a neighbourhood of a local minimum, L-BFGS
is applied for $N_b$ full-batch steps, exploiting the accumulated
curvature information in its inverse Hessian approximation to refine
the solution to high precision.  The strong Wolfe line search of
L-BFGS~\cite{wolfe1969} guarantees sufficient decrease and curvature
conditions at each step, ensuring stable convergence even in the
presence of near-singular Hessian directions.

\subsubsection{GP Surrogate Update and Acquisition}

After the inner loop completes, the final validation loss
$f_t = \mathcal{L}_{val}(\theta^*, \lambda_t)$ is recorded and used
to update the GP posterior.  The GP is refitted by maximising the
marginal likelihood with respect to its kernel hyperparameters
($\sigma_f^2, \ell, \sigma_n^2$) at every trial.  The next candidate
$\lambda_{t+1}$ is obtained by maximising $\mathrm{EI}(\lambda)$ over
the hyperparameter bounds $\mathcal{B}$ using multi-start L-BFGS-B.
This ensures that the acquisition surface is not trapped in local
maxima of EI, which can be highly multimodal.

\subsubsection{Characteristics and Limitations}

Method B is deliberately designed as the strongest practically
competitive baseline.  It combines two of the most effective
single-purpose optimisers (Adam and L-BFGS) with the sample-efficient
BO outer loop, and imposes no particular structural constraints on
the search.  Its key limitation is that each BO trial requires a
complete retraining from scratch, and the BO acquisition function
treats the model as a static black box—there is no mechanism for
continuous HP adaptation during a single training run, nor any
preservation of symplectic geometry.

\subsection{Method C: Unified HHD-ABBO Framework}

\subsubsection{Motivating Principle}

Method C is designed around a single design philosophy: \emph{all
three optimisers should be active within every training epoch}, each
applied at the stage of training where it is most effective.  By
eliminating the separate outer BO loop and instead using the HMC
leapfrog to explore the HP space continuously, Method C recovers the
symplectic guarantees of Method A.  By applying L-BFGS as a
plateau-triggered refinement step, it recovers the second-order
convergence of Method B.  The Adam micro-steps provide the stochastic
gradient signal that drives progress on mini-batches and maintains
adaptivity to data variation.

\subsubsection{Three-Phase Epoch Curriculum}

The unified training curriculum is represented schematically as:
\begin{equation}
\underbrace{\text{Adam}}_{\text{Phase 1:\ fast stochastic descent}}
\;\longrightarrow\;
\underbrace{\text{HMC leapfrog}}_{\text{Phase 2:\ joint }(\theta,\lambda)\text{ co-evolution}}
\;\longrightarrow\;
\underbrace{\text{L-BFGS}}_{\text{Phase 3:\ curvature-based plateau escape}}
\label{eq:curriculum}
\end{equation}

\paragraph{Phase 1 — Adam Micro-Steps}
At the start of each outer epoch, Adam performs $N_{micro}$
gradient steps over randomly sampled mini-batches:
\begin{equation}
\theta \leftarrow \theta - \eta\,\frac{\hat{m}_t}{\sqrt{\hat{v}_t}+\epsilon}
\quad \cite{kingma2015adam}
\end{equation}
This phase maintains rapid stochastic descent and prevents the
Hamiltonian dynamics of Phase 2 from operating on a stale gradient
landscape.

\paragraph{Phase 2 — HMC Leapfrog Proposal}
Once the Adam micro-steps have updated the weights, a single HMC
proposal is generated by running the leapfrog integrator for $L$
steps over the joint phase space $(\theta, \lambda, p_\theta, p_\lambda)$:
\begin{equation}
(\theta,\lambda)
\;\xrightarrow{\;\text{leapfrog }L\text{ steps}\;}
(\theta',\lambda')
\;\xrightarrow{\;\min(1,\,e^{-\Delta H})\;}
\text{accept / reject}
\quad \cite{neal2011}
\end{equation}
The proposal is accepted or rejected by the Metropolis--Hastings rule,
maintaining detailed balance and ergodicity.  Both weights \emph{and}
hyperparameters are updated simultaneously in this symplectic sweep.

\paragraph{Phase 3 — Plateau-Triggered L-BFGS}
After Phase 2, the training loss $\ell_{train}$ is added to a rolling
buffer $\mathcal{B}$ of length $P$.  If the relative loss decrease over
the last $P$ epochs falls below tolerance $\tau$—indicating that the
optimiser has entered a plateau—L-BFGS is activated for $N_b$ full-batch
steps:
\begin{equation}
\theta \leftarrow \theta - \alpha_t\,H_t^{-1}\nabla_\theta\mathcal{L}
\quad \cite{liu1989lbfgs}
\end{equation}
Upon activation, the plateau buffer is reset.  This ensures that
L-BFGS is applied precisely when curvature information is most
valuable—when first-order methods have stalled—and not wastefully
during phases of rapid gradient descent.

\subsubsection{Unique Advantages Over Methods A and B}

Method C offers properties that neither Method A nor Method B can
provide individually:

\begin{itemize}[leftmargin=*]
\item \textbf{Symplectic conservation and detailed balance} (inherited
      from A via HMC~\cite{betancourt2017,leimkuhler2004}): the
      leapfrog integrator preserves a shadow Hamiltonian
      $\tilde{H} = H + \mathcal{O}(\varepsilon^2)$, bounding the
      numerical energy error globally~\cite{hairer2006}.  Method B
      provides no such guarantee.

\item \textbf{Second-order curvature information} (inherited from B via
      L-BFGS~\cite{liu1989lbfgs}): the implicit inverse Hessian
      approximation enables super-linear convergence near stationary
      points.  Stand-alone HMC (Method A) does not exploit curvature
      beyond the mass matrix.

\item \textbf{Intra-epoch efficiency}: all three optimisers operate
      within a single epoch, eliminating the expensive per-trial
      overhead of the outer BO loop.  This yields competitive
      accuracy at a fraction of the wall time of Method B.

\item \textbf{Smooth, continuous HP trajectories}: the HMC dynamics
      generate smooth trajectories $\lambda(t)$ governed by
      Eq.~\eqref{eq:heq3}, providing interpretable and physically
      principled hyperparameter evolution absent from the discrete
      jumps of BO~\cite{zhang2025}.

\item \textbf{Adaptivity to training phase}: the plateau detector
      ensures that L-BFGS is activated only when needed, preventing
      unnecessary computation during rapid-descent phases and
      concentrating second-order refinement at convergence
      boundaries~\cite{rathore2024,kiyani2025}.
\end{itemize}
\section{Algorithms}
This section presents the complete algorithmic descriptions of all
three methods, with detailed annotations explaining the role of each
step.
\subsection{Algorithm 1: Method A — Hamiltonian HMC Training}
Algorithm~\ref{alg:hhd} implements the two-phase HHD procedure.
Phase 1 performs a standard Adam warm-up to place $(\theta, \lambda)$
near a reasonable basin before committing to the Hamiltonian dynamics.
Phase 2 alternates between (i) fresh momentum resampling
$p_\theta \sim \mathcal{N}(0, m_\theta)$,
$p_\lambda \sim \mathcal{N}(0, m_\lambda)$,
(ii) $L$ joint leapfrog steps over the augmented phase space, (iii)
Metropolis--Hastings accept-or-reject using the Hamiltonian error
$\Delta H$, and (iv) Adam micro-steps that continue to refine the
weights within the accepted configuration.  Recording the
Hamiltonian $H = T_\theta + T_\lambda + \mathcal{L}$ at each epoch
enables post-hoc verification of symplectic conservation.
\begin{algorithm}[tb]
\caption{Method A — Hamiltonian HMC Training (HHD)}
\label{alg:hhd}
\textbf{Optimisers:} Adam (Phase 1 warm-up and Phase 2 micro-steps),
Leapfrog integrator + Metropolis--Hastings accept/reject (Phase 2 co-evolution)
\begin{algorithmic}[1]
\Require Initial weights $\theta_0$, initial HPs $\lambda_0$,
         Adam learning rate $\eta$,
         leapfrog step size $\varepsilon$,
         leapfrog sub-steps $L$,
         effective masses $m_\theta$, $m_\lambda$,
         warm-up epochs $N_{wu}$,
         HMC epochs $N_{hmc}$
\State \textbf{Phase 1 — Adam Warm-up}
\For{$t = 1$ to $N_{wu}$}
  \State Compute gradient $g_t \leftarrow \nabla_\theta\mathcal{L}(\theta_{t-1}, \lambda_0)$
  \State Update moments: $m_t$, $v_t$ via Eq.\,\eqref{eq:adam1}--\eqref{eq:adam2}
  \State $\theta_t \leftarrow \theta_{t-1} - \eta\,\hat{m}_t / (\sqrt{\hat{v}_t} + \epsilon)$
         \quad [Eq.\,\eqref{eq:adam3}]
\EndFor
\State Set $\theta \leftarrow \theta_{N_{wu}}$; $\lambda \leftarrow \lambda_0$

\State \textbf{Phase 2 — Joint Hamiltonian Co-evolution}
\For{each epoch $e = 1$ to $N_{hmc}$}
  \State Resample momenta:
         $p_\theta \sim \mathcal{N}(0, m_\theta\,\mathbf{I}_d)$,
         $p_\lambda \sim \mathcal{N}(0, m_\lambda\,\mathbf{I}_k)$
  \State Record initial Hamiltonian:
         $H_{\mathrm{old}} \leftarrow \mathcal{L}(\theta,\lambda)
            + \|p_\theta\|^2/(2m_\theta)
            + \|p_\lambda\|^2/(2m_\lambda)$
         \quad [Eq.\,\eqref{eq:augH}]
  \State Run $L$ leapfrog steps over $(\theta, \lambda, p_\theta, p_\lambda)$:
  \For{$\ell = 1$ to $L$}
    \State $p_\theta \leftarrow p_\theta
              - \tfrac{\varepsilon}{2}\nabla_\theta\mathcal{L}(\theta,\lambda)$
           \quad [Eq.\,\eqref{eq:lf1}]
    \State $p_\lambda \leftarrow p_\lambda
              - \tfrac{\varepsilon}{2}\nabla_\lambda\mathcal{L}(\theta,\lambda)$
    \State $\theta \leftarrow \theta + \varepsilon\,p_\theta/m_\theta$
           \quad [Eq.\,\eqref{eq:lf2}]
    \State $\lambda \leftarrow \lambda + \varepsilon\,p_\lambda/m_\lambda$
    \State $p_\theta \leftarrow p_\theta
              - \tfrac{\varepsilon}{2}\nabla_\theta\mathcal{L}(\theta,\lambda)$
           \quad [Eq.\,\eqref{eq:lf3}]
    \State $p_\lambda \leftarrow p_\lambda
              - \tfrac{\varepsilon}{2}\nabla_\lambda\mathcal{L}(\theta,\lambda)$
  \EndFor
  \State Record proposed Hamiltonian:
         $H_{\mathrm{new}} \leftarrow \mathcal{L}(\theta',\lambda')
            + \|p_\theta'\|^2/(2m_\theta)
            + \|p_\lambda'\|^2/(2m_\lambda)$
  \State $\Delta H \leftarrow H_{\mathrm{new}} - H_{\mathrm{old}}$
  \State Sample $u \sim \mathrm{Uniform}(0, 1)$
  \If{$u < \min(1, e^{-\Delta H})$}
    \State Accept: $(\theta, \lambda) \leftarrow (\theta', \lambda')$
         \hfill \textit{// Metropolis--Hastings acceptance}
  \Else
    \State Reject: restore $(\theta, \lambda)$ from snapshot
  \EndIf
  \State Apply Adam micro-steps to $\theta$ (refine weights within accepted HP)
  \State Snap architecture HPs to nearest valid integer
  \State Record $\mathcal{L}_{train}$, $\mathcal{L}_{val}$,
         $H = T_\theta + T_\lambda + \mathcal{L}$, acceptance rate
\EndFor
\State \textbf{Return} $(\theta^*, \lambda^*)$ minimising $\mathcal{L}_{val}$
\end{algorithmic}
\end{algorithm}

% ─────────────────────────────────────────────────────────────────
\subsection{Algorithm 2: Method B — Hybrid ABBO}

Algorithm~\ref{alg:abbo} implements the outer BO loop combined with
the Adam + L-BFGS inner training curriculum.  The first $N_{rand}$
trials use uniform random sampling to initialise the GP with
sufficiently diverse observations before switching to acquisition-guided
selection.  For each trial, a new model is instantiated with
architecture and optimisation hyperparameters specified by $\lambda_t$,
trained with Adam until stochastic convergence, then refined with
L-BFGS.  The best configuration is tracked globally across all trials.

\begin{algorithm}[tb]
\caption{Method B — Hybrid ABBO}
\label{alg:abbo}
\textbf{Optimisers:} Adam (inner Phase 1), L-BFGS (inner Phase 2),
GP with EI acquisition (outer HP search)
\begin{algorithmic}[1]
\Require GP surrogate with RBF kernel [Eq.\,\eqref{eq:rbf}],
         HP bounds $\mathcal{B}$,
         total BO trials $N_{BO}$,
         random initialisation trials $N_{rand}$,
         Adam epochs per trial $N_a$,
         L-BFGS steps per trial $N_b$

\State Initialise GP with empty observation set $\mathcal{D}_0 = \emptyset$
\State Initialise best result: $f^* \leftarrow \infty$,
       $(\theta^*, \lambda^*) \leftarrow$ None

\For{$t = 1$ to $N_{BO}$}
  \If{$t \le N_{rand}$}
    \State $\lambda_t \sim \mathrm{Uniform}(\mathcal{B})$
           \hfill \textit{// Random exploration phase}
  \Else
    \State $\lambda_t \leftarrow \arg\max_{\lambda \in \mathcal{B}}\,\mathrm{EI}(\lambda)$
           \hfill \textit{// EI acquisition [Eq.\,\eqref{eq:ei}]}
    \State (Maximise via multi-start L-BFGS-B over $\mathcal{B}$)
  \EndIf

  \State Instantiate model $\mathcal{M}(\lambda_t)$ with architecture $\lambda_t$

  \State \textbf{Inner Phase 1 — Adam Training}
  \For{$e = 1$ to $N_a$}
    \State $\theta \leftarrow \theta - \eta\,\hat{m}_e / (\sqrt{\hat{v}_e}+\epsilon)$
           \quad [Eq.\,\eqref{eq:adam3}]
  \EndFor

  \State \textbf{Inner Phase 2 — L-BFGS Refinement}
  \For{$s = 1$ to $N_b$}
    \State Compute full-batch gradient $g_s = \nabla_\theta\mathcal{L}(\theta)$
    \State $\theta \leftarrow \theta - \alpha_s\,H_s^{-1}\,g_s$
           \quad [Eq.\,\eqref{eq:lbfgs}]
           \hfill \textit{// Wolfe line search for $\alpha_s$}
  \EndFor

  \State Evaluate: $f_t \leftarrow \mathcal{L}_{val}(\theta, \lambda_t)$

  \State Update GP: $\mathcal{D}_t \leftarrow \mathcal{D}_{t-1} \cup \{(\lambda_t, f_t)\}$;
         refit GP kernel hyperparameters

  \State Update acquisition: recompute $\mu(\cdot)$, $\sigma(\cdot)$
         over $\mathcal{B}$; set $f^* \leftarrow \min_{s \le t} f_s$

  \If{$f_t < f^*$}
    \State $f^* \leftarrow f_t$;
           $(\theta^*, \lambda^*) \leftarrow (\theta, \lambda_t)$
  \EndIf
\EndFor
\State \textbf{Return} $(\theta^*, \lambda^*)$ achieving best $f^*$
\end{algorithmic}
\end{algorithm}

% ─────────────────────────────────────────────────────────────────
\subsection{Algorithm 3: Method C — Unified HHD-ABBO}

Algorithm~\ref{alg:unified} implements the full three-phase curriculum.
The outer loop iterates over epochs; within each epoch, all three
optimisers are applied in order.  The plateau detector maintains a
rolling buffer of training losses and triggers L-BFGS when the relative
improvement drops below threshold $\tau$ over $P$ consecutive epochs.
This ensures second-order refinement is activated precisely when first-
and second-moment methods have exhausted their progress, and not before.

\begin{algorithm}[tb]
\caption{Method C — Unified HHD-ABBO}
\label{alg:unified}
\textbf{Optimisers:} Adam (Phase 1 warm-up and micro-steps),
Leapfrog HMC (Phase 2 joint co-evolution),
L-BFGS (Phase 3, plateau-triggered)
\begin{algorithmic}[1]
\Require Initial weights $\theta_0$, initial HPs $\lambda_0$,
         Adam learning rate $\eta$,
         leapfrog parameters $\varepsilon$, $L$, $m_\theta$, $m_\lambda$,
         Adam micro-epochs $N_{micro}$,
         L-BFGS steps $N_b$,
         plateau patience $P$,
         plateau tolerance $\tau$,
         warm-up epochs $N_{wu}$,
         total HMC epochs $N_{epochs}$

\State \textbf{Phase 1 Warm-up:} Adam for $N_{wu}$ epochs
       [Eq.\,\eqref{eq:adam1}--\eqref{eq:adam3}]
\State Initialise plateau buffer: $\mathcal{B}_P \leftarrow \emptyset$
\State $(\theta, \lambda) \leftarrow (\theta_{N_{wu}}, \lambda_0)$

\For{epoch $e = 1$ to $N_{epochs}$}

  \State \textbf{[Phase 1]} Adam Micro-steps:
  \For{$j = 1$ to $N_{micro}$}
    \State Sample mini-batch $\mathcal{S}_j$
    \State $\theta \leftarrow \theta - \eta\,\hat{m} / (\sqrt{\hat{v}}+\epsilon)$
           \hfill \textit{// Eq.\ \eqref{eq:adam3} on $\mathcal{S}_j$}
  \EndFor

  \State \textbf{[Phase 2]} HMC Leapfrog Proposal:
  \State Resample momenta
         $p_\theta \sim \mathcal{N}(0, m_\theta\,\mathbf{I})$,
         $p_\lambda \sim \mathcal{N}(0, m_\lambda\,\mathbf{I})$
  \State $H_{\mathrm{old}} \leftarrow \mathcal{L}(\theta,\lambda)
         + \|p_\theta\|^2/(2m_\theta) + \|p_\lambda\|^2/(2m_\lambda)$
  \State Run $L$ joint leapfrog steps:
         $(\theta,\lambda,p_\theta,p_\lambda) \to
          (\theta',\lambda',p_\theta',p_\lambda')$
         \quad [Eq.\,\eqref{eq:lf1}--\eqref{eq:lf3}]
  \State $H_{\mathrm{new}} \leftarrow \mathcal{L}(\theta',\lambda')
         + \|p_\theta'\|^2/(2m_\theta) + \|p_\lambda'\|^2/(2m_\lambda)$
  \State Accept $(\theta,\lambda) \leftarrow (\theta',\lambda')$
         with prob.\ $\min(1,\,e^{-(H_{\mathrm{new}}-H_{\mathrm{old}})})$
  \State Snap architecture HPs to nearest valid integer

  \State Compute training loss $\ell_e \leftarrow \mathcal{L}_{train}(\theta,\lambda)$
  \State Push $\ell_e$ to plateau buffer $\mathcal{B}_P$

  \State \textbf{[Phase 3]} Plateau Detection and L-BFGS Refinement:
  \If{$|\mathcal{B}_P| \ge P$ \textbf{and}
      $(\ell_{e-P} - \ell_e)/\ell_{e-P} < \tau$}
    \For{$s = 1$ to $N_b$}
      \State $\theta \leftarrow \theta - \alpha_s\,H_s^{-1}\nabla_\theta\mathcal{L}(\theta)$
             \quad [Eq.\,\eqref{eq:lbfgs}]
    \EndFor
    \State Reset plateau buffer: $\mathcal{B}_P \leftarrow \emptyset$
  \EndIf

  \State Record $\mathcal{L}_{train}$, $\mathcal{L}_{val}$,
         Hamiltonian $H = T_\theta+T_\lambda+\mathcal{L}$,
         HP trajectory $\lambda_e$, acceptance rate
\EndFor
\State \textbf{Return} $(\theta^*, \lambda^*)$ minimising $\mathcal{L}_{val}$ over all epochs
\end{algorithmic}
\end{algorithm}

\newpage
\section{Harmonic Oscillator Testbed}

\subsection{Physics Motivation and Benchmark Rationale}

The harmonic oscillator serves as the primary physics benchmark for
several compelling reasons.  First, it possesses a closed-form
analytical Hamiltonian, meaning reconstruction errors can be evaluated
exactly against a known ground truth without any approximation or
extrapolation artefacts.  Second, its quadratic energy surface makes the
sensitivity of learned representations to architectural hyperparameters
(depth, width, learning rate) both measurable and interpretable.  Third,
it is the simplest non-trivial example of a Hamiltonian system with
conservation laws, making it the natural starting point for testing
whether an optimiser preserves symplectic structure.  Fourth, its
ubiquity in physics, engineering, and signal processing means that
results on the harmonic oscillator carry significant practical relevance
beyond the benchmark itself.

\subsection{Hamiltonian Formulation}

The classical harmonic oscillator with unit mass $m = 1$ and unit
spring constant $k = 1$ is governed by:
\begin{equation}
H(q,p) = \frac{p^2}{2m} + \frac{1}{2}kq^2
= \frac{p^2}{2} + \frac{q^2}{2}, \quad m = k = 1
\label{eq:ho}
\end{equation}
where $q$ is the generalised position (displacement) and $p$ is the
generalised momentum.  Hamilton's equations for this system produce
elliptic trajectories in phase space:
\begin{equation}
\dot{q} = \frac{\partial H}{\partial p} = p, \qquad
\dot{p} = -\frac{\partial H}{\partial q} = -q
\end{equation}
The energy surface $H(q, p)$ is a paraboloid over the $(q, p)$ plane.
A neural network trained to predict $H(q, p)$ from the pair $(q, p)$ is
therefore learning a smooth, elliptically symmetric surface—a
geometrically clean test of reconstruction fidelity.  The
$p = 0$ slice of the energy surface reduces to the pure potential
$V(q) = \frac{1}{2}q^2$, which serves as a one-dimensional diagnostic
for checking that the network has correctly learned the quadratic shape.

\subsection{Dataset Generation}

The dataset was generated by uniform random sampling of $(q, p)$ pairs
from the phase-space domain $[-4, 4]^2$.  Additive Gaussian noise
$\varepsilon \sim \mathcal{N}(0, 0.05^2)$ was applied to the computed
Hamiltonian values $H(q, p)$ to simulate measurement uncertainty and
prevent overfitting to exact function values.  A total of $N = 800$
samples were generated; these were partitioned into a training set of
640 samples (80\,\%) and a validation set of 160 samples (20\,\%).

Figure~\ref{fig:dataset} presents three complementary views of the
dataset: the ground-truth three-dimensional energy surface
$H(q, p)$, the phase-space contour map illustrating the characteristic
elliptic energy contours, and the 640 training points coloured by their
$H$ values to show coverage of the phase space.

\begin{figure}[tb]
  \centering
  \includegraphics[width=\columnwidth]{plots/fig1_testbed.png}
  \caption{Dataset visualisation.  \textbf{Left:} Ground-truth energy
  surface $H(q,p)$ [Eq.\,\eqref{eq:ho}].
  \textbf{Centre:} Phase-space contour map showing elliptic energy level
  curves.  \textbf{Right:} 640 sampled training points coloured by $H$,
  demonstrating uniform phase-space coverage within $[-4,4]^2$.}
  \label{fig:dataset}
\end{figure}

\subsection{Evaluation Metrics}

Model quality on the harmonic oscillator testbed is assessed using four
complementary metrics:

\textbf{Best Validation MSE:}
\begin{equation}
\mathcal{L}_{val} = \frac{1}{N_{val}}\sum_{i=1}^{N_{val}}
  \bigl(\hat{H}(q_i,p_i;\theta,\lambda) - H(q_i,p_i)\bigr)^2
\label{eq:mse}
\end{equation}

\textbf{Mean Absolute Error (MAE):}  The average absolute deviation
of the reconstructed energy surface from the ground truth, evaluated
over a dense grid in $[-4,4]^2$.

\textbf{Root Mean Squared Error (RMSE):}  The square root of the
mean squared deviation, providing a metric that penalises large errors
more heavily than MAE.

\textbf{Coefficient of Determination ($R^2$):}  A normalised measure
of reconstruction fidelity relative to the variance of the ground-truth
surface; a value of $R^2 = 1$ indicates perfect reconstruction.

Beyond reconstruction quality, two structural properties are assessed:
(i) \textbf{Hamiltonian conservation}: whether the total Hamiltonian
$H = T_\theta + T_\lambda + \mathcal{L}$ [Eq.\,\eqref{eq:augH}] remains
approximately constant over training, and (ii)
\textbf{smoothness of hyperparameter trajectories}: whether $\lambda(t)$
evolves continuously or in discrete jumps.  These structural metrics
distinguish the theoretical quality of each framework independently of
raw numerical accuracy.

% ══════════════════════════════════════════════════════════════════
%  SECTION 7 — EXPERIMENTAL SETUP
% ══════════════════════════════════════════════════════════════════

\section{Experimental Setup}
\label{sec:setup}

\subsection{Common Settings}

All experiments were run on a single CPU core to ensure fair
wall-time comparisons.  The random seed was fixed at 101 across all
three methods for reproducibility.  The MSE loss function
[Eq.\,\eqref{eq:mse}] was used throughout the harmonic oscillator
experiments.  No data augmentation or regularisation beyond dropout
was applied.

\begin{table}[tb]
\centering
\caption{Common Experimental Configuration}
\label{tab:common_setup}
\renewcommand{\arraystretch}{1.2}
\begin{tabular}{ll}
\toprule
Setting & Value \\
\midrule
Dataset (physics)    & Harmonic oscillator, $N = 800$ samples \\
Train / Val split    & $80\,\% / 20\,\%$ (640 / 160 samples) \\
Phase-space domain   & $q, p \in [-4, 4]^2$ \\
Noise standard dev.  & $\sigma = 0.05$ \\
Loss criterion       & MSE: $\mathcal{L} = \frac{1}{N}\sum(\hat{H}_i - H_i)^2$ \\
Random seed          & 101 \\
Device               & CPU \\
\bottomrule
\end{tabular}
\end{table}

\subsection{Method A (HHD) Setup}

Method A employs a 20-epoch Adam warm-up phase followed by 60 epochs
of Hamiltonian co-evolution.  The leapfrog step size $\varepsilon = 0.003$
is deliberately conservative relative to the loss curvature, ensuring
that the numerical Hamiltonian error $\Delta H$ remains small and
acceptance rates stay high.  The effective weight mass $m_\theta = 1.0$
is set to the identity scaling, while the HP mass $m_\lambda = 0.5$
is half the weight mass, providing a balanced exploration rate between
weights and hyperparameters.  The thermostat temperature $T = 1 \times 10^9$
is set to the always-accept regime for optimisation mode.

\begin{table}[tb]
\centering
\caption{Method A (HHD) Configuration}
\label{tab:setup_a}
\renewcommand{\arraystretch}{1.2}
\begin{tabular}{ll}
\toprule
Parameter & Value \\
\midrule
Adam warm-up epochs             & 20 \\
HMC training epochs             & 60 \\
Leapfrog step size $\varepsilon$ & 0.003 \\
Leapfrog sub-steps $L$          & 3 \\
Weight mass $m_\theta$          & 1.0 \\
HP mass $m_\lambda$             & 0.5 \\
Thermostat temperature $T$      & $10^9$ (optimisation mode) \\
HP gradient method              & Central finite differences \\
\bottomrule
\end{tabular}
\end{table}

\subsection{Method B (ABBO) Setup}

Method B runs 15 BO trials in total: 5 trials of random exploration
followed by 10 trials guided by the GP-EI acquisition function
[Eq.\,\eqref{eq:ei}].  Within each trial, Adam is run for 15 epochs
and L-BFGS for 10 full-batch steps.  The GP uses an RBF kernel with
initial length-scale $\ell = 1.0$, signal variance $\sigma_f = 1.0$,
and observation noise $\sigma_n^2 = 10^{-6}$ (essentially noise-free
observations).  The acquisition function is maximised by 5-start
L-BFGS-B.  The exploration bonus is $\xi = 0.01$.

\begin{table}[tb]
\centering
\caption{Method B (ABBO) Configuration}
\label{tab:setup_b}
\renewcommand{\arraystretch}{1.2}
\begin{tabular}{ll}
\toprule
Parameter & Value \\
\midrule
Total BO trials $N_{BO}$          & 15 \\
Random initialisation trials $N_{rand}$ & 5 \\
Adam epochs per trial $N_a$       & 15 \\
L-BFGS steps per trial $N_b$      & 10 \\
GP kernel                         & RBF (squared exponential) \\
GP length-scale $\ell$            & 1.0 (initial) \\
GP signal variance $\sigma_f$     & 1.0 (initial) \\
EI exploration bonus $\xi$        & 0.01 \\
Acquisition maximisation          & 5-start L-BFGS-B \\
\bottomrule
\end{tabular}
\end{table}

\subsection{Method C (Unified HHD-ABBO) Setup}

Method C runs 20 warm-up Adam epochs (cosine-annealed, gradient-clipped)
followed by 60 outer epochs of the three-phase curriculum.  Within each
outer epoch, 3 Adam micro-epochs are applied before the HMC proposal,
ensuring that the weight landscape is refreshed with recent gradient
information before every leapfrog step.  The leapfrog uses 6 sub-steps
with an initial step size of 0.005, adapted via the gradient-norm-aware
controller.  Plateau detection uses a patience window of 4 epochs and a
relative tolerance of $\tau = 5 \times 10^{-5}$; L-BFGS is triggered
with 30 full-batch steps when the condition is met.

\begin{table}[tb]
\centering
\caption{Method C (Unified) Configuration}
\label{tab:setup_c}
\renewcommand{\arraystretch}{1.2}
\begin{tabular}{ll}
\toprule
Parameter & Value \\
\midrule
Adam warm-up epochs              & 20 (cosine-annealed) \\
Outer HMC epochs $N_{epochs}$   & 60 \\
Adam micro-epochs $N_{micro}$   & 3 per outer epoch \\
Leapfrog step size $\varepsilon$ & 0.005 (adaptive) \\
Leapfrog sub-steps $L$          & 6 \\
Weight mass $m_\theta$           & 1.0 \\
HP mass $m_\lambda$              & 0.5 \\
L-BFGS steps $N_b$               & 30 (plateau-triggered) \\
Plateau patience $P$             & 4 epochs \\
Plateau tolerance $\tau$         & $5 \times 10^{-5}$ \\
Gradient clip norm               & 1.0 \\
EMA decay                        & 0.995 \\
\bottomrule
\end{tabular}
\end{table}

\subsection{CNN Classification Setup (MNIST)}

To evaluate the generalisability of the three frameworks beyond
physics benchmarks, each method was applied to train a convolutional
neural network on a balanced subset of the MNIST handwritten digit
classification dataset.  The training set comprised 1000 samples (100
per class) and the validation set 200 samples (20 per class).  The HP
search space was reduced to two dimensions: learning rate
$\eta \in [10^{-4}, 10^{-1}]$ and dropout rate $p_{\mathrm{drop}} \in [0, 0.5]$.
The CNN architecture consisted of two convolutional blocks followed by
two fully-connected layers; depth and width were fixed for a fair
comparison of the optimisation dynamics.

\begin{table}[tb]
\centering
\caption{CNN Classification (MNIST) Configuration}
\label{tab:cnn_setup}
\renewcommand{\arraystretch}{1.2}
\begin{tabular}{ll}
\toprule
Parameter & Value \\
\midrule
Dataset              & MNIST (subset) \\
Training samples     & 5000 \\
Validation samples   & 1000 \\
Architecture         & 2 conv blocks + 2 FC layers \\
HP search space      & $\eta \in [10^{-4}, 10^{-1}]$,
                       $p_{\mathrm{drop}} \in [0, 0.5]$ \\
Loss criterion       & Cross-entropy \\
Metric               & Validation accuracy (\%) \\
Random seed          & 101 \\
Device               & CPU \\
\bottomrule
\end{tabular}
\end{table}

% ══════════════════════════════════════════════════════════════════
%  SECTION 8 — RESULTS
% ══════════════════════════════════════════════════════════════════

\section{Results}

\subsection{Method A: Pure HHD}
Figure~\ref{fig:method_a} provides a detailed analysis of Method A's performance, including loss convergence, Hamiltonian conservation (acceptance rate), energy surface reconstruction, and smooth hyperparameter trajectories.
\begin{figure}[tb]
  \centering
  \includegraphics[width=\columnwidth]{plots/fig2_method_a.png}
  \caption{Method A (Pure HHD) analysis: continuous co-evolution of weights and hyperparameters.}
  \label{fig:method_a}
\end{figure}

\subsection{Method B: Hybrid BO}
Figure~\ref{fig:method_b} displays the results for Method B, showcasing the discrete nature of Bayesian Optimization trials and the highly accurate energy landscape reconstruction achieved at the cost of high computational overhead.
\begin{figure}[tb]
  \centering
  \includegraphics[width=\columnwidth]{plots/fig3_method_b.png}
  \caption{Method B (Hybrid BO) analysis: discrete outer-loop optimization and landscape reconstruction.}
  \label{fig:method_b}
\end{figure}

\subsection{Method C: Unified HHD-ABBO}
Figure~\ref{fig:method_c} illustrates the unified three-phase curriculum of Method C. The approach achieves smooth, continuous hyperparameter evolution while utilizing plateau-triggered L-BFGS refinement to match the precision of Method B in a single training run.
\begin{figure}[tb]
  \centering
  \includegraphics[width=\columnwidth]{plots/fig4_method_c.png}
  \caption{Method C (Unified HHD-ABBO) analysis: three-phase curriculum combining fast descent, symplectic co-evolution, and curvature refinement.}
  \label{fig:method_c}
\end{figure}

\subsection{Standardised Benchmarks: HPOBench, HPOLib, and NAS-Bench-201}

\subsubsection{Benchmark Suite Overview}

To evaluate the generalisation and robustness of the proposed methods
beyond the physics and CNN benchmarks, we conducted comprehensive
benchmarking on 11 standardised tabular datasets spanning three major
HPO benchmark suites:

\begin{itemize}[leftmargin=*]
    \item \textbf{HPOBench~\cite{eggensperger2021,klein2019tabular}:}
          Four classification tasks---\emph{Australian},
          \emph{Blood Transfusion}, \emph{Vehicle}, and
          \emph{Segment}---evaluating discrete hyperparameter search spaces
          for multi-layer perceptrons and support vector machines.
    \item \textbf{HPOLib~\cite{klein2019tabular}:}
          Four regression tasks---\emph{Naval Propulsion},
          \emph{Parkinsons Telemonitoring}, \emph{Protein Structure},
          and \emph{Slice Localisation}---evaluating neural network
          regression hyperparameter configurations.
    \item \textbf{NAS-Bench-201~\cite{dong2020nasbench201}:}
          Architecture search over 15{,}625 cell-based topologies
          evaluated on \emph{CIFAR-10}, \emph{CIFAR-100}, and
          \emph{ImageNet-16-120}~\cite{chrabaszcz2017downsampled}.
\end{itemize}

All three suites employ \emph{tabular lookup tables}: every possible
hyperparameter configuration has been pre-evaluated and stored, so the
benchmark returns the exact validation metric via table lookup.
This paradigm, introduced by Ying \emph{et al.}~\cite{ying2019nasbench101}
and extended by Eggensperger \emph{et al.}~\cite{eggensperger2021},
eliminates training variance, hardware dependence, and random
initialisation effects, ensuring \textbf{100\,\% reproducible results}
that isolate the pure optimisation logic from confounding factors.

\subsubsection{Continuous Relaxation for Discrete Spaces}

Methods A, B, and C were originally designed for continuous
hyperparameter spaces.  To apply them to tabular/discrete benchmarks,
we employ a continuous index relaxation: each discrete parameter
with $N_j$ possible values is mapped to a continuous coordinate
$\lambda_j \in [0, N_j - 1]$, and the nearest integer index is
used for evaluation:
\begin{equation}
\lambda_j^{\mathrm{discrete}} = \mathrm{choices}_j\bigl[
   \mathrm{round}(\mathrm{clip}(\lambda_j, 0, N_j-1))\bigr]
\label{eq:index_relax}
\end{equation}
For GP-based Method B, indices are further normalised to $[0,1]$ by
dividing by the maximum index, ensuring the RBF kernel operates on
a unit-scaled space.  Since no computational graph exists for lookup-table
benchmarks, hyperparameter gradients $\nabla_\lambda\mathcal{L}$
are approximated by one-sided finite differences with perturbation
$\delta = 1.0$ (one index step):
\begin{equation}
\frac{\partial\mathcal{L}}{\partial\lambda_d}
  \approx \frac{\mathcal{L}(\lambda + \delta\,e_d) -
                \mathcal{L}(\lambda)}{\delta}
\label{eq:fd_grad}
\end{equation}

\subsubsection{Experimental Protocol}

We compared five search strategies: Random Search, Optuna
TPE~\cite{akiba2019optuna,bergstra2011tpe}, Method~A (Pure HHD),
Method~B (Hybrid BO), and Method~C (Unified).  Each algorithm was
executed for \textbf{100 trials} across \textbf{5 independent random
seeds} (0--4) on each dataset, yielding a total of
$11 \times 5 \times 100 \times 5 = 27{,}500$ benchmark evaluations.
All seeds are passed to \texttt{numpy.random.RandomState} for
deterministic reproducibility; Optuna TPE uses
\texttt{TPESampler(seed=seed)}.

\begin{table}[tb]
\centering
\caption{Tabular Benchmark Experimental Protocol}
\label{tab:tabular_protocol}
\renewcommand{\arraystretch}{1.2}
\begin{tabular}{ll}
\toprule
Setting & Value \\
\midrule
Benchmark suites     & HPOBench, HPOLib, NAS-Bench-201 \\
Total datasets       & 11 \\
Trials per run       & 100 \\
Seeds                & 5 (0, 1, 2, 3, 4) \\
Warm-up fraction     & 20\,\% (random exploration) \\
HMC leapfrog steps $L$ & 3 \\
HMC step size $\varepsilon$ & 0.5 (adaptive in Method C) \\
GP candidates (Method B) & 1000 random per acquisition \\
\bottomrule
\end{tabular}
\end{table}

\subsubsection{HPOBench Classification Results}

Table~\ref{tab:hpobench_class} reports the mean best validation accuracy
(negated cost) across 5 seeds for each of the four HPOBench classification
datasets.

\begin{table}[tb]
\centering
\caption{HPOBench Classification — Mean Best Accuracy (\%)}
\label{tab:hpobench_class}
\renewcommand{\arraystretch}{1.15}
\begin{tabular}{lccccc}
\toprule
\textbf{Dataset} & \textbf{Random} & \textbf{Optuna} & \textbf{A} & \textbf{B} & \textbf{C} \\
\midrule
Australian        & \textbf{95.72} & 95.17 & 95.61 & 95.61 & 95.61 \\
Blood Transf.     & 72.95 & \textbf{74.72} & 73.72 & 73.24 & 74.60 \\
Vehicle           & 96.89 & 96.50 & \textbf{97.26} & 96.41 & 96.85 \\
Segment           & \textbf{97.23} & 96.94 & 96.94 & 96.94 & \textbf{97.23} \\
\bottomrule
\end{tabular}
\end{table}

On the HPOBench classification suite, Random Search achieves the best
accuracy on Australian (95.72\,\%) and Segment (97.23\,\%, tied with
Method~C) owing to the small, densely covered search space where
uniform exploration is sufficient.  Method~A achieves the best result
on Vehicle (97.26\,\%), demonstrating that HMC exploration can find
configurations that other methods miss in moderately-sized discrete
spaces.  Optuna TPE leads on Blood Transfusion, closely followed by
Method~C.

\subsubsection{HPOLib Regression Results}

Table~\ref{tab:hpolib_reg} reports the mean best validation MSE
across 5 seeds for each HPOLib regression dataset.

\begin{table}[tb]
\centering
\caption{HPOLib Regression — Mean Best Validation MSE}
\label{tab:hpolib_reg}
\renewcommand{\arraystretch}{1.15}
\begin{tabular}{lccccc}
\toprule
\textbf{Dataset} & \textbf{Random} & \textbf{Optuna} & \textbf{A} & \textbf{B} & \textbf{C} \\
\midrule
Naval Prop.     & 1.2e-4 & \textbf{3.1e-5} & 9.9e-5 & 5.1e-5 & 3.4e-5 \\
Parkinsons      & 7.9e-3 & 8.2e-3 & \textbf{7.1e-3} & 1.0e-2 & 7.1e-3 \\
Protein Str.    & 2.8e-1 & \textbf{2.2e-1} & 2.4e-1 & 2.2e-1 & 2.4e-1 \\
Slice Loc.      & 4.6e-4 & 3.2e-4 & 3.8e-4 & 5.2e-4 & \textbf{2.1e-4} \\
\bottomrule
\end{tabular}
\end{table}

The HPOLib results highlight Method~C's versatility.  On Slice
Localisation, Method~C achieves the lowest MSE ($2.1 \times 10^{-4}$),
outperforming all competitors including Optuna TPE.  On Naval
Propulsion, Method~C ranks second ($3.4 \times 10^{-5}$), closely
trailing Optuna ($3.1 \times 10^{-5}$).  Notably, Method~A achieves the
best MSE on Parkinsons ($7.1 \times 10^{-3}$, tied with Method~C),
demonstrating that HMC-based exploration can be competitive on
regression tasks.  Method~B performs poorly on both Parkinsons and
Slice Localisation, suggesting that GP-EI is less effective on these
particular loss surfaces.

\subsubsection{NAS-Bench-201 Architecture Search Results}

Table~\ref{tab:nasbench_results} reports the mean best test accuracy
across 5 seeds on the three NAS-Bench-201 datasets.

\begin{table}[tb]
\centering
\caption{NAS-Bench-201 — Mean Best Test Accuracy (\%)}
\label{tab:nasbench_results}
\renewcommand{\arraystretch}{1.15}
\begin{tabular}{lccccc}
\toprule
\textbf{Dataset} & \textbf{Random} & \textbf{Optuna} & \textbf{A} & \textbf{B} & \textbf{C} \\
\midrule
CIFAR-10       & 91.09 & 91.28 & 91.31 & \textbf{91.60} & 91.42 \\
CIFAR-100      & 72.42 & 73.04 & 72.42 & \textbf{73.52} & 72.60 \\
ImageNet-16    & 45.33 & 46.87 & 45.33 & \textbf{46.90} & 46.53 \\
\bottomrule
\end{tabular}
\end{table}

On NAS-Bench-201, Method~B (GP-EI) consistently achieves the highest
accuracy across all three datasets, which is expected given that the
GP surrogate is particularly well-suited to the smooth, low-dimensional
structure of the NAS-Bench-201 search space (6 categorical operation
choices over DAG edges~\cite{dong2020nasbench201}).  Method~C ranks
second on CIFAR-10 (91.42\,\%) and third on both CIFAR-100 and
ImageNet, while Method~A and Random Search struggle, often obtaining
identical results due to limited exploration in the graph-structured
space.

\subsubsection{Aggregate Rankings and Cross-Benchmark Analysis}

\begin{figure*}[t]
  \centering
  \includegraphics[width=0.48\textwidth]{plots/fig6_hpobench_regret.png}
  \hfill
  \includegraphics[width=0.48\textwidth]{plots/fig6b_hpolib_regret.png}
  \caption{Optimisation trajectories (mean best-so-far objective cost
  $\pm 1$ standard deviation) across 5 seeds on HPOBench classification
  (left) and HPOLib regression (right).}
  \label{fig:tabular_regret}
\end{figure*}

\begin{figure}[tb]
  \centering
  \includegraphics[width=\columnwidth]{plots/fig7_nasbench201_regret.png}
  \caption{Optimisation trajectories on NAS-Bench-201 benchmarks
  (CIFAR-10, CIFAR-100, ImageNet-16-120).}
  \label{fig:nasbench_regret}
\end{figure}

\begin{figure}[tb]
  \centering
  \includegraphics[width=\columnwidth]{plots/fig8_hpobench_summary.png}
  \caption{Cross-benchmark performance summary: mean final best cost
  (top) and per-dataset rank heatmap (bottom) across all 11 datasets.}
  \label{fig:benchmark_summary}
\end{figure}

Figure~\ref{fig:tabular_regret} and Figure~\ref{fig:nasbench_regret}
track the optimisation trajectories (mean best cost vs. trials) across
the three benchmark suites.  Figure~\ref{fig:benchmark_summary} reports
the mean final performance aggregated across all benchmarks alongside a
detailed rank heatmap.

Table~\ref{tab:hpobench_rankings} reports the per-dataset rankings
(1 = best, 5 = worst) recomputed from the actual mean best costs in
\texttt{summary.json}, alongside the average rank across all 11 datasets.

\begin{table}[tb]
\centering
\caption{Optimiser Rankings on Standardised Benchmarks (1 = best)}
\label{tab:hpobench_rankings}
\begin{tabular}{lccccc}
\toprule
\textbf{Benchmark} & \textbf{Random} & \textbf{Optuna} & \textbf{Method A} & \textbf{Method B} & \textbf{Method C} \\
 & \textbf{Search} & \textbf{TPE} & \textbf{(HHD)} & \textbf{(ABBO)} & \textbf{(Unified)} \\
\midrule
Australian          & \textbf{1} & 5 & 2 & 2 & 2 \\
Blood Transfusion   & 5 & \textbf{1} & 3 & 4 & \textbf{2} \\
Vehicle             & 2 & 4 & \textbf{1} & 5 & 3 \\
Segment             & \textbf{1} & 3 & 3 & 3 & \textbf{1} \\
\midrule
Naval Propulsion    & 5 & \textbf{1} & 4 & 3 & \textbf{2} \\
Parkinsons          & 3 & 4 & \textbf{1} & 5 & \textbf{2} \\
Protein Structure   & 5 & \textbf{1} & 4 & 2 & 3 \\
Slice Localisation  & 4 & 2 & 3 & 5 & \textbf{1} \\
\midrule
CIFAR-10            & 5 & 4 & 3 & \textbf{1} & 2 \\
CIFAR-100           & 4 & 2 & 4 & \textbf{1} & 3 \\
ImageNet-16         & 4 & 2 & 4 & \textbf{1} & 3 \\
\midrule
\textbf{Avg Rank}   & \textbf{3.55} & \textbf{2.64} & \textbf{2.91} & \textbf{2.91} & \textbf{2.18} \\
\bottomrule
\end{tabular}
\end{table}

Method~C (Unified) achieves the \textbf{best average rank of 2.18}
across all 11 tabular benchmarks, outperforming Optuna TPE (2.64),
Method~A and Method~B (both 2.91), and Random Search (3.55).
Notably, Method~C obtains a rank of 1 or 2 on 7 out of 11 datasets,
demonstrating that the three-phase curriculum generalises effectively
to discrete hyperparameter spaces via continuous index relaxation.

The per-suite analysis reveals distinct performance patterns.
On \textbf{HPOBench classification}, Method~C and Random Search both
perform well, while Optuna TPE is penalised on Australian (rank 5)
where the search space is small enough that random exploration
suffices.  On \textbf{HPOLib regression}, Method~C achieves ranks
1--2 on three of four datasets, with Optuna TPE excelling on
Naval Propulsion and Protein Structure.  On \textbf{NAS-Bench-201},
Method~B (GP-EI) dominates due to its strength on smooth,
low-dimensional graph-structured search spaces~\cite{dong2020nasbench201},
achieving rank 1 on all three architecture search datasets.

The adaptive step-size controller in Method~C, which dynamically adjusts
the leapfrog $\varepsilon$ based on acceptance rate, proves
critical for discrete spaces where the optimal exploration rate varies
across datasets.  Without adaptation (Method~A, fixed $\varepsilon = 0.5$),
the HMC explorer either over-rejects (too many proposals land far from
valid integer configurations) or under-explores, resulting in a
substantially worse average rank (2.91 vs. 2.18).

\subsubsection{Reproducibility}

All tabular benchmark results are fully reproducible.  The
\texttt{simple-hpo-bench}~\cite{klein2019tabular} library provides
deterministic lookup tables, and all random number generators are
seeded explicitly via \texttt{numpy.random.RandomState(seed)}.
Individual trajectory files are stored per-seed per-dataset in
JSON format, enabling independent verification.  The complete benchmark
can be reproduced with:
\begin{verbatim}
pip install simple-hpo-bench numpy optuna
            scikit-learn scipy
python hpobench_benchmark.py
python plot_hpobench.py
\end{verbatim}

\subsection{Comparative Analysis and Multi-Metric Comparison}
Figure~\ref{fig:comparative} and Table~\ref{tab:results} present a comprehensive multi-metric comparison across all three methods. Method C achieves the strongest profile across all quantitative metrics, surpassing Method B's reconstruction accuracy on the harmonic oscillator (MSE $0.00327$ vs $0.09896$) while retaining the symplectic and structural properties of Method A.
\begin{figure}[tb]
  \centering
  \includegraphics[width=\columnwidth]{plots/fig5_comparative.png}
  \caption{Comparative Overview: convergence comparison, reconstruction quality, and CNN classification benchmarks.}
  \label{fig:comparative}
\end{figure}

\begin{table}[tb]
\centering
\caption{Quantitative Results --- Harmonic Oscillator Benchmark}
\label{tab:results}
\renewcommand{\arraystretch}{1.25}
\begin{tabular}{lccc}
\toprule
\textbf{Metric}     & \textbf{A: HMC} & \textbf{B: Hybrid} & \textbf{C: Unified} \\
\midrule
Best Val.\ MSE      & 0.15128 & 0.09896   & \textbf{0.00327} \\
Landscape MAE       & 0.6990  & 0.1050    & \textbf{0.0177}  \\
Landscape RMSE      & 0.7863  & 0.1394    & \textbf{0.0248}  \\
$R^2$               & 0.9498  & 0.9984    & \textbf{0.9999}  \\
Wall time (s)       & \textbf{17.8}   & 147.1   & 81.4    \\
H conservation      & \textbf{Yes}   & No      & \textbf{Yes}    \\
Smooth HP traj.\ & \textbf{Yes}   & No      & \textbf{Yes}    \\
Detailed balance    & \textbf{Yes}   & No      & \textbf{Yes}    \\
2nd-order curv.\    & No      & \textbf{Yes}    & \textbf{Yes}    \\
\bottomrule
\end{tabular}
\end{table}

\begin{table}[tb]
\centering
\caption{Quantitative Results --- CNN Classification on MNIST}
\label{tab:cnn_results}
\renewcommand{\arraystretch}{1.25}
\begin{tabular}{lccc}
\toprule
\textbf{Metric}     & \textbf{A: HMC} & \textbf{B: Hybrid} & \textbf{C: Unified} \\
\midrule
Best Val.\ Acc (\%) & \textbf{97.80} & 97.40 & 97.70 \\
Final Val.\ Acc (\%) & \textbf{97.80} & 97.40 & 97.70 \\
Wall time (s)       & \textbf{122.7}  & 698.7  & 214.8  \\
Final LR $\eta$     & $1.1\!\times\!10^{-3}$
                    & $6.2\!\times\!10^{-4}$
                    & $7.9\!\times\!10^{-4}$ \\
Final Dropout $p_{\mathrm{drop}}$ & 0.21 & 0.22 & 0.24 \\
\bottomrule
\end{tabular}
\end{table}


% ══════════════════════════════════════════════════════════════════
%  SECTION 9 — DISCUSSION
% ══════════════════════════════════════════════════════════════════


\section{Discussion}

\subsection{Theoretical Properties of the Unified Framework}

The most significant theoretical contribution of Method C is that it
is the \emph{only} framework among the three that simultaneously
provides all four desirable properties: symplectic conservation,
detailed balance, second-order curvature information, and efficient
intra-epoch compute.

\begin{proposition}[Per-Epoch Computational Cost of Method~C]
\label{prop:complexity}
Let $d$ denote the total number of network weight parameters,
$k$ the dimensionality of the hyperparameter space,
$L$ the number of leapfrog steps per HMC proposal,
$m$ the L-BFGS history size,
$n$ the mini-batch size,
and $t$ the number of completed BO trials.  Then:
\begin{enumerate}
\item \textbf{Method~A and Method~C} (per epoch):
\[
  \mathcal{O}\bigl(n \cdot d + L \cdot (n \cdot d + k \cdot d) + n \cdot d\bigr)
  = \mathcal{O}\bigl(L \cdot d \cdot (n + k)\bigr),
\]
consisting of Adam micro-step $\mathcal{O}(nd)$, $L$ leapfrog steps
each requiring forward/backward passes of cost $\mathcal{O}(nd)$ and
finite-difference HP gradients of cost $\mathcal{O}(kd)$, and a
Metropolis accept/reject check of cost $\mathcal{O}(nd)$.  The L-BFGS
plateau step (when triggered) adds $\mathcal{O}(md)$ per step over
$N_b$ steps, giving $\mathcal{O}(N_b m d)$ amortised over the plateau
detection window.
\item \textbf{Method~B} (per BO trial):
\[
  \mathcal{O}\bigl(N_a \cdot n \cdot d + N_b \cdot m \cdot d + t^3\bigr),
\]
consisting of Adam training of cost $\mathcal{O}(N_a n d)$, L-BFGS
refinement of cost $\mathcal{O}(N_b m d)$, and GP posterior update of
cost $\mathcal{O}(t^3)$ scaling \emph{cubically} with the number of
completed trials.
\end{enumerate}
The \emph{key distinction} is that Method~B's GP inference cost grows
as $\mathcal{O}(t^3)$, whereas Method~C eliminates this term entirely,
with each epoch costing $\mathcal{O}(L \cdot d \cdot (n + k))$
independently of the number of epochs elapsed.
\end{proposition}

\begin{proof}[Proof sketch]
Adam's per-step cost follows directly from its $\mathcal{O}(d)$
moment update and parameter step.  Each leapfrog half-step requires
one forward and one backward pass for $\nabla_\theta\mathcal{L}$
(cost $\mathcal{O}(nd)$) plus a central finite-difference evaluation
for each of the $k$ hyperparameters (cost $\mathcal{O}(k \cdot nd)$
with two forward passes per HP).  The GP posterior update requires
solving a $t \times t$ linear system via Cholesky decomposition, which
costs $\mathcal{O}(t^3)$.  L-BFGS applies the two-loop recursion
using $m$ vector pairs of dimension $d$, giving $\mathcal{O}(md)$
per step.
\end{proof}

The \textbf{symplectic conservation} property, inherited from the
leapfrog integrator, means that the shadow Hamiltonian
$\tilde{H} = H + \mathcal{O}(\varepsilon^2)$ is preserved globally
throughout training~\cite{hairer2006}.  This is not merely an academic
guarantee: in practice, it ensures that the Metropolis--Hastings
acceptance rate remains high and that the HP trajectory does not drift
arbitrarily far from the energy minimum.  The Hamiltonian conservation
plot (Figure~\ref{fig:method_c}) confirms this empirically.

The \textbf{detailed balance} condition guarantees that the Markov
chain defined by the HMC proposals is ergodic and converges to the
correct stationary distribution over the joint $(\theta, \lambda)$
space~\cite{betancourt2017}.  This provides a principled probabilistic
interpretation of the joint training process that is absent from both
BO (which does not define a stationary distribution over weights) and
standard gradient descent.

The \textbf{second-order curvature} information from L-BFGS provides
super-linear convergence near stationary points.  The implicit inverse
Hessian approximation identifies the principal curvature directions of
the loss surface and scales the gradient step accordingly, resolving
ill-conditioned loss landscapes that stymie Adam's diagonal
preconditioner~\cite{liu1989lbfgs}.  The plateau detection mechanism
ensures that this expensive computation is amortised over multiple
epochs, with L-BFGS activated only when its benefit exceeds its cost.

\subsection{Interpretability of Hyperparameter Trajectories}

A key practical advantage of Methods A and C that is difficult to
quantify in a single scalar metric is the \emph{interpretability} of
their hyperparameter trajectories.  Because the HMC dynamics are
governed by Hamilton's Eq.\,\eqref{eq:heq3}—a physically motivated
differential equation—the trajectory $\lambda(t)$ has a natural
interpretation as the path of a particle in the hyperparameter
potential $\mathcal{L}(\theta, \lambda)$.

This interpretability has several practical consequences.  First, it
allows practitioners to visualise and diagnose the HP optimisation
process: a trajectory that oscillates around a fixed point indicates
a well-behaved convergence, while a drifting trajectory suggests that
the loss landscape has a shallow slope in that HP dimension.  Second,
the smooth trajectories provide implicit regularisation against
overfitting to a particular hyperparameter configuration, as the inertia
of the momenta prevents abrupt switches.  Third, the continuous evolution
of HPs within a single training run eliminates the statistical
discontinuities introduced by BO's discrete restarts, providing a more
stable final configuration.

\subsection{Wall-Time Efficiency and Practical Deployability}

From a practical deployment perspective, the wall-time results are
decisive.  Method B's 147-second runtime for the harmonic oscillator
benchmark---compared to 17.8 seconds for A and 81.4 seconds for C---reflects
the fundamental overhead of running 15 complete training loops, each
with its own model initialisation, Adam training, and L-BFGS refinement,
followed by GP inference and acquisition maximisation.  As the size of
the model or dataset grows, this overhead scales multiplicatively with
the number of BO trials, making Method B progressively less attractive
relative to the intra-epoch curriculum of Method C.

Method C achieves a $1.8\times$ speedup over Method B (81.4\,s vs
147.1\,s) on the harmonic oscillator while simultaneously achieving
\textbf{superior reconstruction quality} (MSE $0.00327$ vs $0.09896$,
$R^2 = 0.9999$ vs $0.9984$).  On the CNN benchmark, Method~C achieves
$97.70\,\%$ accuracy in $214.8$\,s---a $3.3\times$ speedup over
Method~B's $698.7$\,s.  On the 11 standardised tabular benchmarks,
Method~C attains the best average rank ($2.18$), outperforming
Optuna TPE ($2.64$) and both individual frameworks.  This combination
of superior or matched accuracy and lower compute cost makes the Unified
framework the preferred choice across all evaluated settings.

\subsection{When to Choose Each Method}

\begin{table}[tb]
\centering
\caption{Method Selection Guide}
\label{tab:guide}
\renewcommand{\arraystretch}{1.35}
\begin{tabular}{p{1.8cm}p{5.5cm}}
\toprule
\textbf{Method} & \textbf{Preferred scenarios} \\
\midrule
\textbf{A — HHD} & Theoretical guarantees are paramount (posterior
sampling, uncertainty quantification, physical interpretability,
energy conservation).  Maximum speed with moderate accuracy is
acceptable.  Bayesian inference over weights is desired~\cite{betancourt2017}. \\
\midrule
\textbf{B — ABBO} & Strong GP-surrogate-based search on smooth,
low-dimensional spaces (e.g., NAS-Bench-201).  Wall time is not a
constraint.  Parallel trial evaluation is available to amortise BO
overhead~\cite{snoek2012}.  No requirement for continuous HP
trajectories or symplectic properties. \\
\midrule
\textbf{C — Unified} & Best all-round choice across accuracy,
theoretical guarantees, computational efficiency, and HP
interpretability.  Achieves the strongest reconstruction accuracy
on energy landscapes and the best average rank on tabular HPO
benchmarks.  Physics-informed or structure-preserving training
is required.  Resource-constrained settings where both fast
convergence and high accuracy are needed~\cite{kiyani2025}. \\
\bottomrule
\end{tabular}
\end{table}

\subsection{Limitations and Failure Modes}

Method A's HMC component is sensitive to the leapfrog step size
$\varepsilon$ and the mass ratio $m_\lambda / m_\theta$.  A step size
that is too large results in large $\Delta H$ and low acceptance rates;
a step size that is too small produces highly correlated proposals that
explore the phase space slowly.  Automatic step-size adaptation via the
No-U-Turn Sampler (NUTS)~\cite{betancourt2017} or dual averaging would
address this limitation in future work.

Method B is inherently sequential in its standard implementation:
each BO trial requires the previous trial's results before the next
acquisition can be computed.  Parallel BO variants (e.g., using the
batch Expected Improvement acquisition) could significantly reduce the
wall-time disadvantage.

Method C's plateau detector requires careful tuning of patience $P$
and tolerance $\tau$.  A poorly calibrated detector may trigger L-BFGS
too early (wasting compute on non-convergent phases) or too late
(missing opportunities to escape plateaux).  Adaptive plateau
detection based on loss curvature rather than relative change would
improve robustness~\cite{rathore2024}.

% ══════════════════════════════════════════════════════════════════
%  SECTION 10 — CONCLUSION
% ══════════════════════════════════════════════════════════════════

\section{Conclusion}

This paper presented a rigorous comparative study of three complementary
hyperparameter optimisation frameworks for neural energy landscape
learning, evaluated across physics benchmarks, CNN classification,
and 11 standardised tabular HPO benchmarks.

\textbf{Method A (HHD)} established that joint weight--hyperparameter
co-evolution via symplectic leapfrog integration is feasible and
produces smooth, physically interpretable HP trajectories with full
Hamiltonian conservation.  Its wall time of 17.8 seconds on the harmonic
oscillator benchmark confirms that the Hamiltonian dynamics do not
impose prohibitive computational overhead.

\textbf{Method B (ABBO)} confirmed that the combination of Adam,
L-BFGS, and GP-guided Bayesian Optimisation provides a strong practical
baseline, achieving competitive reconstruction quality (MAE $0.1050$,
$R^2 = 0.9984$) on the harmonic oscillator and the best results on
NAS-Bench-201 architecture search tasks.  Its primary limitation---the
$147.1$-second runtime driven by 15 complete training runs---restricts
its deployability in compute-constrained scenarios.

\textbf{Method C (Unified HHD-ABBO)} is the principal novel contribution
of this work.  By integrating Adam micro-steps, HMC leapfrog proposals,
and plateau-triggered L-BFGS refinement into a single three-phase
epoch curriculum, Method C inherits symplectic conservation and detailed
balance from Method A, and second-order curvature from Method B, while
eliminating the expensive outer BO loop entirely.  The result is the
\textbf{best reconstruction accuracy} on the harmonic oscillator
(MSE $= 0.00327$, $R^2 = 0.9999$), outperforming Method~B by
$30\times$ in error magnitude, $97.70\,\%$ CNN validation accuracy
in $3.3\times$ less wall time than Method~B, and the \textbf{best
average rank (2.18)} across all 11 standardised tabular benchmarks,
surpassing Optuna TPE (2.64), Random Search (3.55), and both individual
frameworks (2.91).

The multi-benchmark evaluation---spanning the harmonic oscillator
testbed, MNIST CNN classification, and 11 datasets from HPOBench,
HPOLib, and NAS-Bench-201---comprehensively demonstrates that the
Unified framework generalises across physics, classification, regression,
and architecture search tasks.  The multi-metric comparison, spanning
quantitative accuracy metrics (MAE, RMSE, $R^2$, MSE), qualitative
structural properties (Hamiltonian conservation, smooth HP trajectories,
detailed balance, second-order curvature), computational efficiency
(wall time), and cross-benchmark rankings, provides a thorough
characterisation of each framework's strengths.

\subsection*{Future Directions}

Several natural extensions of this work are identified:

\begin{enumerate}[leftmargin=*, label=\arabic*.]
\item \textbf{Automatic step-size adaptation:}  Implementing the
      dual-averaging scheme of NUTS~\cite{betancourt2017} within
      the leapfrog integrator to remove the need for manual tuning
      of $\varepsilon$ and improve robustness.

\item \textbf{Parallelised HMC proposals:}  Generating multiple
      HMC proposals in parallel per epoch to increase the effective
      exploration rate of the HP space without increasing wall time
      proportionally.

\item \textbf{Extension to physics-informed neural networks
      (PINNs):}  Applying the Unified framework to PINNs trained on
      partial differential equations~\cite{krishnapriyan2021,wang2023,kiyani2025},
      where the HP landscape is substantially more complex and
      symplectic conservation of the learned dynamics is a hard
      physical constraint.

\item \textbf{Adaptive plateau detection:}  Replacing the fixed
      patience-and-tolerance criterion with a curvature-adaptive
      trigger based on the Hessian spectrum, ensuring that L-BFGS
      is activated precisely when second-order information is most
      beneficial~\cite{rathore2024}.

\item \textbf{Scalability to large-scale models:}  Evaluating the
      framework on transformer-scale models where the high-dimensional
      weight space $d \gg 10^6$ requires distributed leapfrog
      integration and memory-efficient L-BFGS approximations.
\end{enumerate}

% ══════════════════════════════════════════════════════════════════
%  REFERENCES
% ══════════════════════════════════════════════════════════════════

\bibliographystyle{siamplain}
\begin{thebibliography}{99}

\bibitem{kingma2015adam}
D.~P. Kingma and J.~Ba,
``Adam: A Method for Stochastic Optimization,''
in \textit{Proc. Int. Conf. Learning Representations (ICLR)}, 2015.

\bibitem{liu1989lbfgs}
D.~C. Liu and J.~Nocedal,
``On the Limited Memory BFGS Method for Large Scale Optimization,''
\textit{Mathematical Programming}, vol.~45, no.~1--3, pp.~503--528, 1989.

\bibitem{wolfe1969}
P.~Wolfe,
``Convergence Conditions for Ascent Methods,''
\textit{SIAM Review}, vol.~11, no.~2, pp.~226--235, 1969.

\bibitem{byrd1995}
R.~H. Byrd, P.~Lu, J.~Nocedal, and C.~Zhu,
``A Limited Memory Algorithm for Bound Constrained Optimization,''
\textit{SIAM J. Scientific Computing}, vol.~16, no.~5,
pp.~1190--1208, 1995.

\bibitem{mockus1975}
J.~Mockus,
``On Bayesian Methods for Seeking the Extremum,''
in \textit{Optimization Techniques IFIP Technical Conference}, 1975.

\bibitem{snoek2012}
J.~Snoek, H.~Larochelle, and R.~P. Adams,
``Practical Bayesian Optimization of Machine Learning Algorithms,''
in \textit{Advances in Neural Information Processing Systems (NeurIPS)},
vol.~25, 2012.

\bibitem{rasmussen2006}
C.~E. Rasmussen and C.~K.~I. Williams,
\textit{Gaussian Processes for Machine Learning}.
MIT Press, 2006.

\bibitem{duane1987}
S.~Duane, A.~D. Kennedy, B.~J. Pendleton, and D.~Roweth,
``Hybrid Monte Carlo,''
\textit{Physics Letters B}, vol.~195, no.~2, pp.~216--222, 1987.

\bibitem{neal2011}
R.~M. Neal,
``MCMC Using Hamiltonian Dynamics,''
in \textit{Handbook of Markov Chain Monte Carlo}, CRC Press, 2011,
ch.~5, pp.~113--162.

\bibitem{betancourt2017}
M.~Betancourt,
``A Conceptual Introduction to Hamiltonian Monte Carlo,''
\textit{arXiv preprint arXiv:1701.02434}, 2017.

\bibitem{leimkuhler2004}
B.~Leimkuhler and S.~Reich,
\textit{Simulating Hamiltonian Dynamics}.
Cambridge University Press, 2004.

\bibitem{hairer2006}
E.~Hairer, C.~Lubich, and G.~Wanner,
\textit{Geometric Numerical Integration: Structure-Preserving Algorithms
for Ordinary Differential Equations}, 2nd ed.
Springer, 2006.

\bibitem{bergstra2012}
J.~Bergstra and Y.~Bengio,
``Random Search for Hyper-Parameter Optimization,''
\textit{Journal of Machine Learning Research}, vol.~13,
pp.~281--305, 2012.

\bibitem{lecun2015}
Y.~LeCun, Y.~Bengio, and G.~Hinton,
``Deep Learning,''
\textit{Nature}, vol.~521, no.~7553, pp.~436--444, 2015.

\bibitem{greydanus2019}
S.~Greydanus, M.~Dzamba, and J.~Yosinski,
``Hamiltonian Neural Networks,''
in \textit{Advances in Neural Information Processing Systems (NeurIPS)},
vol.~32, 2019.

\bibitem{zhang2025}
Z.~Zhang, S.~Bai, Y.~Liang, and Z.~Bao,
``Neural Networks Under Hamiltonian Constraints: A Comprehensive Review,''
\textit{IEEE Access}, 2025.

\bibitem{goldstein2002}
H.~Goldstein, C.~Poole, and J.~Safko,
\textit{Classical Mechanics}, 3rd ed.
Addison-Wesley, 2002.

\bibitem{raiaan2024}
M.~A.~K. Raiaan \textit{et al.},
``A Systematic Review of Hyperparameter Optimization Techniques in CNNs,''
\textit{Decision Analytics Journal}, vol.~11, p.~100470, 2024.

\bibitem{frazier2018}
P.~I. Frazier,
``A Tutorial on Bayesian Optimization,''
\textit{arXiv preprint arXiv:1807.02811}, 2018.

\bibitem{rathore2024}
P.~Rathore, W.~Lei, Z.~Frangella, L.~Lu, and M.~Udell,
``Challenges in Training PINNs: A Loss Landscape Perspective,''
in \textit{Proc. 41st ICML}, PMLR 235, 2024.

\bibitem{krishnapriyan2021}
A.~Krishnapriyan \textit{et al.},
``Characterizing Possible Failure Modes in Physics-Informed Neural Networks,''
in \textit{NeurIPS}, vol.~34, pp.~26548--26560, 2021.

\bibitem{wang2023}
S.~Wang, S.~Sankaran, H.~Wang, and P.~Perdikaris,
``An Expert's Guide to Training Physics-Informed Neural Networks,''
\textit{arXiv preprint arXiv:2308.08468}, 2023.

\bibitem{kiyani2025}
E.~Kiyani \textit{et al.},
``Optimizing the Optimizer for PINNs and Kolmogorov-Arnold Networks,''
\textit{arXiv preprint arXiv:2501.16371}, 2025.

\bibitem{jin2022}
M.~Jin \textit{et al.},
``Hyperparameter Tuning of ANNs for Well Production Estimation,''
\textit{ACS Omega}, vol.~7, no.~28, pp.~24145--24156, 2022.

\bibitem{chen2024}
W.~Chen, A.~A. Howard, and P.~Stinis,
``Self-Adaptive Weights Based on Balanced Residual Decay Rate for PINNs,''
\textit{arXiv preprint arXiv:2404.18055}, 2024.

\bibitem{newhouse2024}
L.~Newhouse,
\textit{Mathematical Techniques for Tuning Hyperparameters of RNNs}.
Research Monograph, 2024.

\bibitem{wang2021}
S.~Wang, Y.~Teng, and P.~Perdikaris,
``Understanding and Mitigating Gradient Pathologies in Physics-Informed
Neural Networks,''
\textit{SIAM J. Scientific Computing}, 2021.

\bibitem{wong2025}
J.~C.~Wong \textit{et al.},
``Evolutionary Optimization of PINNs: Evo-PINN Frontiers,''
\textit{arXiv preprint arXiv:2501.06572}, 2025.

\bibitem{eggensperger2021}
K.~Eggensperger \textit{et al.},
``HPOBench: A Collection of Reproducible Multi-Fidelity Benchmark Problems for HPO,''
in \textit{NeurIPS Datasets and Benchmarks}, 2021.

\bibitem{paszke2019}
A.~Paszke \textit{et al.},
``PyTorch: An Imperative Style, High-Performance Deep Learning Library,''
in \textit{NeurIPS}, vol.~32, 2019.

\bibitem{virtanen2020}
P.~Virtanen \textit{et al.},
``SciPy 1.0: Fundamental Algorithms for Scientific Computing in Python,''
\textit{Nature Methods}, vol.~17, pp.~261--272, 2020.

% ── New references (HPO Benchmarks) ──

\bibitem{klein2019tabular}
A.~Klein and F.~Hutter,
``Tabular Benchmarks for Joint Architecture and Hyperparameter Optimization,''
\textit{arXiv preprint arXiv:1905.04970}, 2019.

\bibitem{ying2019nasbench101}
C.~Ying, A.~Klein, E.~Christiansen, E.~Real, K.~Murphy, and F.~Hutter,
``NAS-Bench-101: Towards Reproducible Neural Architecture Search,''
in \textit{Proc. 36th ICML}, PMLR 97, pp.~7105--7114, 2019.

\bibitem{dong2020nasbench201}
X.~Dong and Y.~Yang,
``NAS-Bench-201: Extending the Scope of Reproducible Neural Architecture
Search,'' in \textit{Proc. ICLR}, 2020.

\bibitem{chrabaszcz2017downsampled}
P.~Chrabaszcz, I.~Loshchilov, and F.~Hutter,
``A Downsampled Variant of ImageNet as an Alternative to the CIFAR Datasets,''
\textit{arXiv preprint arXiv:1707.08819}, 2017.

% ── New references (TPE / Optuna) ──

\bibitem{bergstra2011tpe}
J.~Bergstra, R.~Bardenet, Y.~Bengio, and B.~K\'{e}gl,
``Algorithms for Hyper-Parameter Optimization,''
in \textit{Advances in Neural Information Processing Systems (NeurIPS)},
vol.~24, pp.~2546--2554, 2011.

\bibitem{akiba2019optuna}
T.~Akiba, S.~Sano, T.~Yanase, T.~Ohta, and M.~Koyama,
``Optuna: A Next-Generation Hyperparameter Optimization Framework,''
in \textit{Proc. 25th ACM SIGKDD}, pp.~2623--2631, 2019.

\bibitem{watanabe2023tpe}
S.~Watanabe,
``Tree-Structured Parzen Estimator: Understanding Its Theory and Analyzing
a Connection to Bayesian Optimization,''
\textit{arXiv preprint arXiv:2304.11127}, 2023.

% ── New references (Bayesian Optimization) ──

\bibitem{shahriari2016bo}
B.~Shahriari, K.~Swersky, Z.~Wang, R.~P. Adams, and N.~de~Freitas,
``Taking the Human Out of the Loop: A Review of Bayesian Optimization,''
\textit{Proceedings of the IEEE}, vol.~104, no.~1, pp.~148--175, 2016.

\bibitem{hutter2011smac}
F.~Hutter, H.~H. Hoos, and K.~Leyton-Brown,
``Sequential Model-Based Optimization for General Algorithm Configuration,''
in \textit{Proc. 5th Int. Conf. Learning and Intelligent Optimization
(LION)}, Springer LNCS 6683, pp.~507--523, 2011.

\bibitem{li2018hyperband}
L.~Li, K.~Jamieson, G.~DeSalvo, A.~Rostamizadeh, and A.~Talwalkar,
``Hyperband: A Novel Bandit-Based Approach to Hyperparameter Optimization,''
\textit{Journal of Machine Learning Research}, vol.~18, no.~185,
pp.~1--52, 2018.

\bibitem{falkner2018bohb}
S.~Falkner, A.~Klein, and F.~Hutter,
``BOHB: Robust and Efficient Hyperparameter Optimization at Scale,''
in \textit{Proc. 35th ICML}, PMLR 80, pp.~1437--1446, 2018.

% ── New references (Neural Architecture Search) ──

\bibitem{zoph2017nas}
B.~Zoph and Q.~V. Le,
``Neural Architecture Search with Reinforcement Learning,''
in \textit{Proc. 5th ICLR}, 2017.

\bibitem{elsken2019nas}
T.~Elsken, J.~H. Metzen, and F.~Hutter,
``Neural Architecture Search: A Survey,''
\textit{Journal of Machine Learning Research}, vol.~20, no.~55,
pp.~1--21, 2019.

\bibitem{liu2019darts}
H.~Liu, K.~Simonyan, and Y.~Yang,
``DARTS: Differentiable Architecture Search,''
in \textit{Proc. 7th ICLR}, 2019.

\bibitem{real2019aging}
E.~Real, A.~Aggarwal, Y.~Huang, and Q.~V. Le,
``Regularized Evolution for Image Classifier Architecture Search,''
in \textit{Proc. 33rd AAAI Conf. Artificial Intelligence}, vol.~33,
pp.~4780--4789, 2019.

% ── New references (HMC Theory) ──

\bibitem{hoffman2014nuts}
M.~D. Hoffman and A.~Gelman,
``The No-U-Turn Sampler: Adaptively Setting Path Lengths in Hamiltonian
Monte Carlo,''
\textit{Journal of Machine Learning Research}, vol.~15, no.~47,
pp.~1593--1623, 2014.

\bibitem{girolami2011riemann}
M.~Girolami and B.~Calderhead,
``Riemann Manifold Langevin and Hamiltonian Monte Carlo Methods,''
\textit{Journal of the Royal Statistical Society: Series B},
vol.~73, no.~2, pp.~123--214, 2011.

% ── New references (Optimization Theory) ──

\bibitem{nocedal2006numerical}
J.~Nocedal and S.~J. Wright,
\textit{Numerical Optimization}, 2nd ed.
Springer, 2006.

\bibitem{ruder2016sgd}
S.~Ruder,
``An Overview of Gradient Descent Optimization Algorithms,''
\textit{arXiv preprint arXiv:1609.04747}, 2016.

\bibitem{loshchilov2019adamw}
I.~Loshchilov and F.~Hutter,
``Decoupled Weight Decay Regularization,''
in \textit{Proc. 7th ICLR}, 2019.

% ── New references (Physics-Informed ML) ──

\bibitem{raissi2019pinns}
M.~Raissi, P.~Perdikaris, and G.~E. Karniadakis,
``Physics-Informed Neural Networks: A Deep Learning Framework for Solving
Forward and Inverse Problems Involving Nonlinear Partial Differential
Equations,'' \textit{Journal of Computational Physics}, vol.~378,
pp.~686--707, 2019.

\bibitem{cranmer2020lagrangian}
M.~Cranmer, S.~Greydanus, S.~Hoyer, P.~Battaglia, D.~Spergel, and S.~Ho,
``Lagrangian Neural Networks,''
\textit{arXiv preprint arXiv:2003.04630}, 2020.

\bibitem{toth2020hamiltonian}
P.~Toth, D.~J. Rezende, A.~Jaegle, S.~Racani\`{e}re, A.~Botev,
and I.~Higgins,
``Hamiltonian Generative Networks,''
in \textit{Proc. 8th ICLR}, 2020.

% ── New references (Reproducibility) ──

\bibitem{bouthillier2019reproducibility}
X.~Bouthillier, P.~Delaunay, M.~Bronzi \textit{et al.},
``Unreproducible Research is Reproducible,''
in \textit{Proc. 36th ICML}, PMLR 97, pp.~725--734, 2019.

\bibitem{lindauer2020hpo}
M.~Lindauer and F.~Hutter,
``Best Practices for Scientific Research on Neural Architecture Search,''
\textit{Journal of Machine Learning Research}, vol.~21, no.~243,
pp.~1--18, 2020.

\\bibitem{lorraine2020}
J.~Lorraine, P.~Vicol, and D.~Duvenaud,
``Optimizing Millions of Hyperparameters by Implicit Differentiation,''
in \textit{Proc. Int. Conf. Artificial Intelligence and Statistics (AISTATS)},
2020, pp.~1--9.

\\bibitem{franceschi2018}
L.~Franceschi, P.~Frasconi, S.~Salzo, R.~Grazzi, and M.~Pontil,
``Bilevel Programming for Hyperparameter Optimization and Meta-Learning,''
in \textit{Proc. Int. Conf. Machine Learning (ICML)},
2018, pp.~1568--1577.

\\bibitem{finn2017maml}
C.~Finn, P.~Abbeel, and S.~Levine,
``Model-Agnostic Meta-Learning for Fast Adaptation of Deep Networks,''
in \textit{Proc. Int. Conf. Machine Learning (ICML)},
2017, pp.~1126--1135.

\\bibitem{franca2020}
G.~Fran\c{c}a, D.~Robinson, and R.~Vidal,
``Conformal Symplectic and Relativistic Optimization,''
\textit{Journal of Statistical Mechanics: Theory and Experiment},
vol.~2020, no.~12, p.~124008, 2020.

\\bibitem{real2019regularized}
E.~Real, A.~Aggarwal, Y.~Huang, and Q.~V.~Le,
``Regularized Evolution for Image Classifier Architecture Search,''
in \textit{Proc. AAAI Conf. Artificial Intelligence},
2019, pp.~4780--4789.

\\bibitem{li2020asha}
L.~Li, K.~Jamieson, A.~Rostamizadeh, E.~Gonina, J.~Ben-Tzur,
M.~Hardt, B.~Recht, and A.~Talwalkar,
``A System for Massively Parallel Hyperparameter Tuning,''
\textit{Proc. Machine Learning and Systems (MLSys)},
vol.~2, pp.~230--246, 2020.

\\bibitem{cranmer2020}
M.~Cranmer, S.~Greydanus, S.~Hoyer, P.~Battaglia, D.~Spergel, and S.~Ho,
``Lagrangian Neural Networks,''
\textit{arXiv preprint arXiv:2003.04630}, 2020.

\end{thebibliography}

\end{document}